\documentclass[11pt]{article}

\usepackage[T1]{fontenc}
\usepackage{lmodern}
\usepackage[a4paper,margin=30mm]{geometry}
\usepackage{amsmath,amssymb,amsthm,mathtools}
\usepackage{microtype}
\usepackage{enumitem}
\usepackage{booktabs,longtable,array}
\usepackage{xcolor}
\usepackage{hyperref}
\usepackage[nameinlink,noabbrev]{cleveref}
\usepackage{listings}

\hypersetup{
  colorlinks=true,
  linkcolor=blue!55!black,
  citecolor=green!35!black,
  urlcolor=blue!65!black,
  pdftitle={Pontryagin Duality for Locally Compact Abelian Groups},
  pdfauthor={Proof blueprint for Mathlib autoformalization}
}

\lstdefinelanguage{Lean}{
  morekeywords={def,theorem,lemma,instance,noncomputable,section,namespace,end,where,by,
    variable,open,scoped,deriving,inferInstance,exact,refine,fun,Type,Prop},
  sensitive=true,
  morecomment=[l]{--},
  morecomment=[s]{/-}{-/},
  morestring=[b]"
}
\lstset{
  language=Lean,
  basicstyle=\ttfamily\small,
  keywordstyle=\color{blue!55!black},
  commentstyle=\color{green!35!black},
  stringstyle=\color{orange!55!black},
  columns=fullflexible,
  keepspaces=true,
  breaklines=true,
  frame=single,
  rulecolor=\color{black!20},
  xleftmargin=0.5em,
  xrightmargin=0.5em
}

\newtheorem{theorem}{Theorem}[section]
\newtheorem{proposition}[theorem]{Proposition}
\newtheorem{lemma}[theorem]{Lemma}
\newtheorem{corollary}[theorem]{Corollary}
\newtheorem{claim}[theorem]{Claim}
\theoremstyle{definition}
\newtheorem{definition}[theorem]{Definition}
\newtheorem{construction}[theorem]{Construction}
\newtheorem{example}[theorem]{Example}
\theoremstyle{remark}
\newtheorem{remark}[theorem]{Remark}

\newcommand{\T}{\mathbb T}
\newcommand{\C}{\mathbb C}
\newcommand{\R}{\mathbb R}
\newcommand{\Z}{\mathbb Z}
\newcommand{\N}{\mathbb N}
\newcommand{\one}{\mathbf 1}
\newcommand{\e}{\mathrm e}
\newcommand{\id}{\mathrm{id}}
\newcommand{\ev}{\operatorname{ev}}
\newcommand{\supp}{\operatorname{supp}}
\newcommand{\Span}{\operatorname{span}}
\newcommand{\Hom}{\operatorname{Hom}}
\newcommand{\im}{\operatorname{im}}
\newcommand{\argu}{\operatorname{arg}}
\newcommand{\Cc}{C_{\mathrm c}}
\newcommand{\Cb}{C_{\mathrm b}}
\newcommand{\Czero}{C_0}
\newcommand{\Lp}[1]{L^{#1}}
\newcommand{\dual}[1]{\widehat{#1}}
\newcommand{\bidual}[1]{\widehat{\widehat{#1}}}
\newcommand{\polar}[1]{#1^{\perp}}
\newcommand{\inv}{\vphantom{f}^{\vee}}
\newcommand{\abs}[1]{\lvert #1\rvert}
\newcommand{\norm}[1]{\lVert #1\rVert}
\newcommand{\set}[1]{\left\{#1\right\}}
\newcommand{\angles}[1]{\left\langle #1\right\rangle}
\newcommand{\mathlib}{\textsf{Mathlib}}
\newcommand{\Circle}{\texttt{Circle}}
\newcommand{\PD}{\texttt{PontryaginDual}}

\newenvironment{formalization}
  {\begin{quote}\small\noindent\textbf{Formalization note.}\ }
  {\end{quote}}

\title{Pontryagin Duality for Locally Compact Abelian Groups\\
\large A self-contained proof blueprint aligned with \mathlib}
\author{Prepared for autoformalization}
\date{Mathlib interface checked against v4.32.0 (July 2026)}

\begin{document}
\maketitle

\begin{abstract}
Let $G$ be a Hausdorff locally compact abelian topological group and let
\[
  \dual G=\Hom_{\mathrm{cts}}(G,\T)
\]
carry the compact--open topology.  We give a complete proof that the evaluation map
\[
  \alpha_G:G\longrightarrow\bidual G,
  \qquad \alpha_G(x)(\chi)=\chi(x),
\]
is an isomorphism of topological groups.  The proof is analytic: compact--open topology,
Haar measure, the commutative Banach algebra $\Lp1(G)$, positive-type functions, the GNS
construction, Bochner's theorem, Fourier inversion, and Plancherel are developed in a
dependency order suited to formalization.  We then prove annihilator duality, duality for
closed subgroups and quotients, extension of characters, exactness, compact/discrete
duality, naturality, and the standard examples.

The target definition is exactly \mathlib's multiplicative definition
\texttt{PontryaginDual A := A ->\_t* Circle}.  Generic results that should be imported
from the library---Haar existence, Riesz--Markov, Fubini--Tonelli, Banach--Alaoglu,
Krein--Milman, and Stone--Weierstrass---are isolated in an explicit foundation interface.
Every argument specific to Pontryagin duality is proved below; no form of biduality,
character extension, or an LCA structure theorem is used as an input.
\end{abstract}

\tableofcontents

\section{Formalization contract and conventions}
\label{sec:contract}

\subsection{The mathematical category}

Throughout the main proof, an \emph{LCA group} means a type $G$ equipped with a
commutative group law and a Hausdorff locally compact topology for which multiplication
and inversion are continuous.  We use multiplicative notation, including for $\R$ and
$\Z$ when they occur as abstract groups.  The circle is
\[
  \T=\set{z\in\C:\abs z=1}
\]
with multiplication and its usual topology.  A \emph{character} of $G$ is a continuous
group homomorphism $G\to\T$.

No countability, metrizability, compact generation, or second-countability hypothesis is
made.  Nets are used whenever sequential reasoning would introduce such a hypothesis.
All measures are regular Borel (Radon) measures.  Equalities of $\Lp p$ functions are
equalities almost everywhere; an equality between continuous functions obtained almost
everywhere is upgraded to pointwise equality because Haar measure has full support.

\subsection{The exact \mathlib interface}

In \mathlib~v4.32.0 the relevant source file is
\path{Mathlib/Topology/Algebra/PontryaginDual.lean}.  Its central definition is:

\begin{lstlisting}
def PontryaginDual := A ->_t* Circle
deriving TopologicalSpace
\end{lstlisting}

Here \texttt{A ->\_t* B} is \path{ContinuousMonoidHom A B}; its topology is induced
from the compact--open topology on continuous maps.  The file already provides:

\begin{enumerate}[label=(M\arabic*)]
\item a commutative group structure on \texttt{PontryaginDual A};
\item Hausdorff and topological-group instances;
\item local compactness of \texttt{PontryaginDual G} when $G$ is a locally compact
  topological group;
\item compactness of the dual of a discrete group and discreteness of the dual of a
  compact group;
\item the contravariant map
\begin{lstlisting}
def PontryaginDual.map (f : A ->_t* B) :
    PontryaginDual B ->_t* PontryaginDual A :=
  f.compLeft Circle
\end{lstlisting}
  together with \texttt{map\_apply}, \texttt{map\_one}, \texttt{map\_comp},
  \texttt{map\_mul}, and \texttt{mapHom}.
\end{enumerate}

The intended new definition and theorem have the following shape.  Names are proposed,
not asserted to exist already.

\begin{lstlisting}
noncomputable def PontryaginDual.eval
    (G : Type*) [CommGroup G] [TopologicalSpace G]
    [IsTopologicalGroup G] [LocallyCompactSpace G] :
    G ->_t* PontryaginDual (PontryaginDual G) := ...

theorem PontryaginDual.eval_bijective
    [T2Space G] : Function.Bijective (PontryaginDual.eval G) := ...

noncomputable def PontryaginDual.toDoubleDual
    [T2Space G] :
    ContinuousMulEquiv G (PontryaginDual (PontryaginDual G)) := ...
\end{lstlisting}

The defining simp lemma should be
\begin{lstlisting}
@[simp] theorem eval_apply (x : G) (chi : PontryaginDual G) :
    PontryaginDual.eval G x chi = chi x := rfl
\end{lstlisting}
and naturality should be exposed in a form that rewrites without category theory:
\begin{lstlisting}
theorem map_map_eval (f : G ->_t* H) :
    ContinuousMonoidHom.comp
      (PontryaginDual.map (PontryaginDual.map f))
      (PontryaginDual.eval G)
      = ContinuousMonoidHom.comp (PontryaginDual.eval H) f := ...
\end{lstlisting}

\subsection{Foundation interface}
\label{subsec:foundations}

Calling the account ``self-contained'' means that every Pontryagin-specific construction
and inference is present.  Re-proving the foundations of measure theory and functional
analysis inside the same document would obscure the formalization boundary.  The proof
uses precisely the following generic results, all stated here so that no hidden
mathematical hypothesis is introduced.

\begin{longtable}{>{\raggedright\arraybackslash}p{0.10\textwidth}
                  >{\raggedright\arraybackslash}p{0.82\textwidth}}
\toprule
Label & Imported statement \\
\midrule
\endhead
(F1) & A locally compact Hausdorff space is regular.  If $x\in U$ with $U$ open, there is
$u\in\Cc(X,[0,1])$ with $u(x)=1$ and $\supp u\subseteq U$.  Compact subsets have compact
neighbourhoods. \\
(F2) & Every locally compact Hausdorff group has a nonzero left Haar measure, unique up
to a positive scalar.  Haar measure is Radon, finite on compact sets, positive on nonempty
open sets, and has full support.  On an abelian group it is also right invariant and
invariant under inversion. \\
(F3) & Tonelli, Fubini, dominated convergence, monotone convergence, change of variables
under a measure-preserving homeomorphism, density of $\Cc(G)$ in $\Lp p(G)$ for
$1\leq p<\infty$, H\"older, and the Riesz representation theorem for Hilbert spaces. \\
(F4) & Riesz--Markov: positive linear functionals on $\Cc(X)$ that are locally bounded
are represented by unique Radon measures.  Finite complex Radon measures are the dual of
$\Czero(X)$; their total variation is regular. \\
(F5) & Tychonoff's theorem and the Arzel\`a--Ascoli theorem in compact--open form:
a closed, pointwise relatively compact, equicontinuous family of continuous maps is
compact on a locally compact domain. \\
(F6) & Banach--Alaoglu for dual balls; Krein--Milman for compact convex subsets of a
locally convex space; Stone--Weierstrass for self-adjoint, point-separating subalgebras of
$\Czero(X)$ that vanish nowhere. \\
(F7) & The standard construction of quotient topological groups by closed normal
subgroups, and completion of a positive semidefinite sesquilinear space to a Hilbert
space. \\
(F8) & Unitization and the spectral-radius bound for Banach algebras, and the spectral
theorem (equivalently, continuous functional calculus and spectral projections) for
bounded normal operators on a complex Hilbert space. \\
\bottomrule
\end{longtable}

The current library contains the principal ingredients of (F1)--(F8), though some glue
lemmas in exactly the stated generality may need to be added.  In particular, the proof
must not replace nets by sequences or invoke an $\Lp1$ duality theorem with an unstated
$\sigma$-finiteness assumption.

\subsection{Fourier conventions}
\label{subsec:conventions}

Fix a Haar measure $\mu_G$ on $G$.  For $a\in G$ set
\[
  (L_a f)(x)=f(a^{-1}x),\qquad
  f^*(x)=\overline{f(x^{-1})},\qquad
  (f*g)(x)=\int_G f(y)g(y^{-1}x)\,d\mu_G(y).
\]
For $f\in\Lp1(G)$ define
\begin{equation}
  \widehat f(\chi)=\int_G f(x)\overline{\chi(x)}\,d\mu_G(x).
  \label{eq:fourier-convention}
\end{equation}
For a finite complex measure $\sigma$ on $\dual G$, define its inverse
Fourier--Stieltjes transform by
\begin{equation}
  \check\sigma(x)=\int_{\dual G}\chi(x)\,d\sigma(\chi).
  \label{eq:FS-convention}
\end{equation}
Thus the desired inversion formula is
\begin{equation}
  f(x)=\int_{\dual G}\widehat f(\chi)\chi(x)\,d\mu_{\dual G}(\chi).
  \label{eq:desired-inversion}
\end{equation}
These conventions give
\begin{align}
  \widehat{f*g}&=\widehat f\,\widehat g,
  &\widehat{f^*}&=\overline{\widehat f},
  &\widehat{L_a f}(\chi)&=\overline{\chi(a)}\,\widehat f(\chi).
  \label{eq:fourier-identities}
\end{align}
All later conjugations and inverses are fixed by these three identities.

\section{The circle and the compact--open character group}
\label{sec:compact-open}

\subsection{Small arcs and the no-small-subgroups lemma}

For $0<r<\pi$ put
\[
  A_r=\set{z\in\T:\abs{\argu z}<r},
\]
where $\argu z\in(-\pi,\pi]$ is the principal argument.

\begin{lemma}[No small subgroups]
\label{lem:nss-circle}
If $z^n\in A_{\pi/2}$ for every $n\in\Z$, then $z=1$.  Consequently the only subgroup
of $\T$ contained in $A_{\pi/2}$ is the trivial subgroup.
\end{lemma}

\begin{proof}
Write $z=e^{i\theta}$ with $\theta\in(-\pi,\pi]$.  If $\theta\neq0$, put
$t=\abs\theta\in(0,\pi]$ and take
$n=\lceil\pi/(2t)\rceil$.  If $t>\pi/2$, then $n=1$; if $t\leq\pi/2$, then
\[
  \frac\pi2\leq nt<\frac\pi2+t\leq\pi.
\]
In either case $nt\in[\pi/2,\pi]$.  Therefore the principal argument of $z^n$ has
absolute value at least $\pi/2$, contradicting $z^n\in A_{\pi/2}$.  Hence $\theta=0$.
In \mathlib this is packaged as
\path{Circle.eq_one_of_forall_pow_mem_centeredArc_pi_div_two}.
\end{proof}

The following quantitative form is what compactness and equicontinuity arguments use.

\begin{lemma}[Power control]
\label{lem:power-control}
Let $W$ be a neighbourhood of $1$ in $\T$.  There is $m\geq1$ such that, for every
$z\in\T$,
\[
  z,z^2,\ldots,z^m\in \overline{A_{\pi/4}}
  \quad\Longrightarrow\quad z\in W.
\]
\end{lemma}

\begin{proof}
If not, for every $m$ choose $z_m\notin W$ whose first $m$ powers lie in the closed arc.
Compactness of $\T$ gives a convergent subnet $z_m\to z\notin W$ after first replacing
$W$ by a neighbourhood with closed closure contained in the original $W$.  For each
fixed $n$, eventually $m\geq n$, so $z_m^n\in\overline{A_{\pi/4}}$ and hence
$z^n\in\overline{A_{\pi/4}}\subset A_{\pi/2}$.  The same holds for negative powers.
\Cref{lem:nss-circle} gives $z=1$, contradicting $z\notin W$.
\end{proof}

\begin{formalization}
The local-compactness instance in the current source uses the explicit basis
$A_{\pi/2^{n+1}}$ and a one-step square-root lemma instead of
\cref{lem:power-control}.  Reusing that exact basis will substantially shorten Lean proofs.
\end{formalization}

\subsection{The compact--open topology}

For compact $K\subseteq G$ and open $U\subseteq\T$ write
\[
  [K,U]=\set{f\in C(G,\T):f(K)\subseteq U}.
\]
The compact--open topology on $C(G,\T)$ is generated by these sets, and $\dual G$ has
the induced topology.  At the trivial character it is enough to take sets $[K,U]$ with
$U$ a neighbourhood of $1$.

\begin{lemma}[Neighbourhood basis]
\label{lem:dual-neighbourhood-basis}
The sets
\[
  W(K,U)=\set{\chi\in\dual G:\chi(K)\subseteq U},
\]
where $K$ is compact and $U$ is a neighbourhood of $1$ in $\T$, form a neighbourhood
basis at $1\in\dual G$.
\end{lemma}

\begin{proof}
A basic neighbourhood of the constant map $1$ is a finite intersection
$\bigcap_j[K_j,U_j]$.  Put $K=\bigcup_jK_j$ and choose a neighbourhood
$U\subseteq\bigcap_jU_j$ of $1$.  Then $W(K,U)$ is contained in the intersection.
Conversely each $W(K,U)$ is the intersection of $\dual G$ with a subbasic compact--open
set, hence is open in $\dual G$.
\end{proof}

\begin{proposition}[Joint evaluation]
\label{prop:joint-evaluation}
If $G$ is locally compact Hausdorff, then
\[
  \ev:\dual G\times G\longrightarrow\T,\qquad(\chi,x)\longmapsto\chi(x)
\]
is continuous.
\end{proposition}

\begin{proof}
Fix $(\chi_0,x_0)$ and a neighbourhood $O$ of $\chi_0(x_0)$.  Choose a symmetric
neighbourhood $V$ of $1$ with $VV\chi_0(x_0)\subseteq O$.  By continuity of $\chi_0$,
there is a neighbourhood $N$ of $x_0$ on which
$\chi_0(x)\chi_0(x_0)^{-1}\in V$.  By local compactness shrink $N$ so that its closure
$K$ is compact.  If $\chi\chi_0^{-1}\in W(K,V)$ and $x\in N$, then
\[
  \chi(x)\chi_0(x_0)^{-1}
  =(\chi\chi_0^{-1})(x)\,
    \chi_0(x)\chi_0(x_0)^{-1}\in VV.
\]
Thus $W(K,V)\chi_0\times N$ is mapped into $O$.
\end{proof}

\begin{formalization}
This is already abstracted by the \texttt{ContinuousEval} instance for continuous
monoid homomorphisms under \texttt{LocallyCompactPair}.  Currying uses the existing
lemma \texttt{continuous\_of\_continuous\_uncurry} in the
\texttt{ContinuousMonoidHom} namespace.
\end{formalization}

\subsection{The dual is again locally compact abelian}

\begin{theorem}
\label{thm:dual-lca}
If $G$ is an LCA group, then $\dual G$, with pointwise multiplication and the
compact--open topology, is an LCA group.
\end{theorem}

\begin{proof}
Pointwise multiplication and inversion preserve characters.  Their continuity follows
from the induced compact--open topology and continuity of the corresponding pointwise
operations on $C(G,\T)$.  Since $\T$ is Hausdorff and two distinct characters differ at
some point, the compact--open function space, and hence $\dual G$, is Hausdorff.

It remains to prove local compactness.  Choose a compact neighbourhood $K$ of $1$ in
$G$ and an arc $A_r$ whose closure is contained in $A_{\pi/4}$.  Let
\[
  E=\set{\chi\in\dual G:\chi(K)\subseteq\overline{A_r}}.
\]
This family is closed in $\dual G$ and pointwise relatively compact because $\T$ is
compact.  It is equicontinuous at $1$.  Indeed, given a neighbourhood $W$ of $1$ in
$\T$, take $m$ from \cref{lem:power-control}.  Choose a neighbourhood $V$ of $1$ in
$G$ such that $V^j\subseteq K$ for $1\leq j\leq m$.  For $v\in V$ and $\chi\in E$,
all powers $\chi(v)^j=\chi(v^j)$ lie in $\overline{A_r}$, so $\chi(v)\in W$.
Translation and the homomorphism identity turn equicontinuity at $1$ into
equicontinuity everywhere.  Arzel\`a--Ascoli now makes $E$ compact.  Finally
$W(K,A_r)$ is an open neighbourhood of $1$ whose closure is contained in $E$.
Thus $\dual G$ has a compact neighbourhood of its identity.
\end{proof}

\begin{formalization}
The theorem is already an instance.  Its proof in v4.32.0 invokes
\texttt{ContinuousMonoidHom.locallyCompactSpace\_of\_hasBasis} with the centred arcs
$A_{\pi/2^{n+1}}$; it ultimately uses Arzel\`a--Ascoli.  The argument above records the
mathematical invariants that must remain visible if that implementation is refactored.
\end{formalization}

\begin{proposition}[Compact/discrete directions]
\label{prop:compact-discrete-basic}
\begin{enumerate}[label=(\alph*)]
\item If $G$ is discrete, then $\dual G$ is compact.
\item If $G$ is compact, then $\dual G$ is discrete.
\end{enumerate}
\end{proposition}

\begin{proof}
If $G$ is discrete, every homomorphism $G\to\T$ is continuous.  The set of
homomorphisms is closed in the compact product $\T^G$, because the equations
$f(1)=1$ and $f(xy)=f(x)f(y)$ are closed conditions.  Compact subsets of a discrete
space are finite, so the product and compact--open topologies agree on this closed set.

If $G$ is compact, $W(G,A_{\pi/2})$ contains only the trivial character: the image of
any character in this set is a subgroup of $\T$ contained in $A_{\pi/2}$, hence is
trivial by \cref{lem:nss-circle}.  The singleton $\{1\}$ is therefore open, and
translations show that every singleton is open.
\end{proof}

\section{Functoriality and the evaluation morphism}
\label{sec:evaluation}

For a continuous homomorphism $f:G\to H$, define
\[
  \dual f:\dual H\longrightarrow\dual G,
  \qquad (\dual f)(\chi)=\chi\circ f.
\]
Continuity follows from continuity of precomposition in the compact--open topology.
Direct evaluation proves
\[
  \dual{\id_G}=\id_{\dual G},\qquad
  \dual{(g\circ f)}=(\dual f)\circ(\dual g).
\]
This is exactly \texttt{PontryaginDual.map} and \texttt{map\_comp}.

\begin{proposition}[Evaluation morphism]
\label{prop:evaluation-map}
For an LCA group $G$, the formula
\[
  \alpha_G(x)(\chi)=\chi(x)
\]
defines a continuous homomorphism $\alpha_G:G\to\bidual G$.
\end{proposition}

\begin{proof}
For fixed $x$, evaluation is multiplicative in $\chi$, so $\alpha_G(x)$ is a group
homomorphism.  It is continuous by continuity of point evaluation on the compact--open
space.  The identity
$\alpha_G(xy)(\chi)=\chi(xy)=\chi(x)\chi(y)$ makes $\alpha_G$ a homomorphism.
Finally the uncurried map $(x,\chi)\mapsto\chi(x)$ is continuous by
\cref{prop:joint-evaluation}; the compact--open exponential law, in the precise form
\texttt{continuous\_of\_continuous\_uncurry}, says that its curry $\alpha_G$ is
continuous.
\end{proof}

\begin{proposition}[Naturality and the triangle identity]
\label{prop:naturality}
For every continuous homomorphism $f:G\to H$,
\begin{equation}
  \bidual f\circ\alpha_G=\alpha_H\circ f.
  \label{eq:naturality}
\end{equation}
Moreover
\begin{equation}
  (\dual{\alpha_G})\circ\alpha_{\dual G}=\id_{\dual G}.
  \label{eq:triangle}
\end{equation}
\end{proposition}

\begin{proof}
For $x\in G$ and $\psi\in\dual H$,
\[
 (\bidual f(\alpha_G(x)))(\psi)
 =\alpha_G(x)(\psi\circ f)=\psi(f(x))=\alpha_H(f(x))(\psi).
\]
For $\chi\in\dual G$ and $x\in G$,
\[
 ((\dual{\alpha_G})(\alpha_{\dual G}(\chi)))(x)
 =\alpha_{\dual G}(\chi)(\alpha_G(x))=\alpha_G(x)(\chi)=\chi(x).
\]
Extensionality of continuous homomorphisms proves both identities.
\end{proof}

\section{Haar measure, convolution, and approximate identities}
\label{sec:haar-convolution}

Fix an LCA group $G$ and a nonzero Haar measure $\mu_G$.  Commutativity implies that
left translation, right translation, and inversion preserve $\mu_G$.

\begin{lemma}[Normalized compactly supported bumps]
\label{lem:normalized-bump}
For every identity neighbourhood $U$ there is $h_U\in\Cc(G)$ such that
\[
  h_U\geq0,\qquad \supp h_U\subseteq U,\qquad
  \int_Gh_U\,d\mu_G=1.
\]
The support may be required to lie in any prescribed relatively compact identity
neighbourhood contained in $U$.
\end{lemma}

\begin{proof}
Choose a nonempty open set $V$ with compact closure contained in $U$.  By (F1), choose
a nonzero $u\in\Cc(G,[0,1])$ supported in $U$ and positive at a point of $V$.
Continuity makes $u$ positive on a nonempty open set, so full support of Haar measure
gives $0<\int u<\infty$.  Put $h_U=u/(\int u)$.
\end{proof}

Order the identity neighbourhoods by reverse inclusion and make one choice $h_U$ for
each $U$.  This is the approximate identity used below.

\begin{lemma}[Continuity of translation]
\label{lem:translation-continuity}
For $1\leq p<\infty$ and $f\in\Lp p(G)$,
\[
  \norm{L_af-f}_p\longrightarrow0\qquad(a\to1).
\]
Consequently $a\mapsto L_af$ is continuous from $G$ to $\Lp p(G)$.
\end{lemma}

\begin{proof}
First take $f\in\Cc(G)$.  For $a$ in a fixed compact neighbourhood of $1$, all
$L_af-f$ have support in one compact set.  Uniform continuity of $f$ on a slightly
larger compact set gives $\norm{L_af-f}_\infty\to0$, and multiplication by the finite
measure of the common support gives convergence in $\Lp p$.  For general $f$, choose
$g\in\Cc(G)$ with $\norm{f-g}_p<\varepsilon$ and use the isometry of translation:
\[
 \norm{L_af-f}_p
 \leq 2\norm{f-g}_p+\norm{L_ag-g}_p.
\]
Continuity at an arbitrary $a_0$ follows after translating by $a_0^{-1}$.
\end{proof}

\begin{proposition}[Approximate identity]
\label{prop:approximate-identity}
For $f\in\Lp p(G)$, $1\leq p<\infty$,
\[
  h_U*f\longrightarrow f\quad\text{in }\Lp p(G)
\]
as $U$ shrinks to $\{1\}$.  If $f$ is continuous, convergence is uniform on each
compact set.  In particular, for every compact $K\subseteq\dual G$,
\begin{equation}
  \widehat{h_U}\longrightarrow1
  \quad\text{uniformly on }K.
  \label{eq:approx-fourier}
\end{equation}
\end{proposition}

\begin{proof}
Minkowski's integral inequality and \cref{lem:translation-continuity} give
\[
 \norm{h_U*f-f}_p
 \leq\int_Uh_U(y)\norm{L_yf-f}_p\,d\mu_G(y),
\]
which is at most $\varepsilon$ once $U$ is small enough.  If $f$ is continuous, the same
estimate with the supremum over a fixed compact set uses uniform continuity on a compact
enlargement of that set.

For \eqref{eq:approx-fourier}, joint evaluation
$(y,\chi)\mapsto\chi(y)$ is uniformly continuous on $C\times K$, where $C$ is a fixed
compact neighbourhood containing all sufficiently small $U$.  Hence
\[
 \sup_{\chi\in K}\abs{\widehat{h_U}(\chi)-1}
 \leq \sup_{\substack{y\in U\\\chi\in K}}
       \abs{\overline{\chi(y)}-1}\longrightarrow0.
\]
\end{proof}

\begin{proposition}[The convolution Banach $*$-algebra]
\label{prop:l1-banach-algebra}
The formulas in \cref{subsec:conventions} make $\Lp1(G)$ a commutative Banach
$*$-algebra, possibly without a unit.  More precisely,
\begin{align*}
 \norm{f*g}_1&\leq\norm f_1\norm g_1,&
 (f*g)*h&=f*(g*h),& f*g&=g*f,\\
 (f*g)^*&=g^**f^*,& \norm{f^*}_1&=\norm f_1.
\end{align*}
Moreover, if $f,g\in\Lp2(G)$, then $f*g$ has a continuous bounded representative and
\begin{equation}
  \norm{f*g}_\infty\leq\norm f_2\norm g_2.
  \label{eq:l2-convolution-bound}
\end{equation}
\end{proposition}

\begin{proof}
Tonelli gives
\[
 \int_G\abs{(f*g)(x)}\,dx
 \leq\int_G\int_G\abs{f(y)}\abs{g(y^{-1}x)}\,dy\,dx
 =\norm f_1\norm g_1.
\]
Fubini and the substitutions permitted by Haar invariance prove associativity,
commutativity, and the involution identity.  Completeness is completeness of $\Lp1$.
For $f,g\in\Lp2$, Cauchy--Schwarz gives \eqref{eq:l2-convolution-bound}.  If $x_i\to x$,
the difference is bounded by $\norm{L_{x_i^{-1}x}f-f}_2\norm g_2$, which tends to zero
by \cref{lem:translation-continuity}.
\end{proof}

\begin{lemma}[Positive-type kernels from convolution]
\label{lem:convolution-positive-type}
If $g\in\Lp2(G)$, then $g*g^*$ has a continuous bounded representative and is of
positive type in the sense of \cref{def:positive-type} below.  If $g\in\Cc(G)$, then
\[
  \supp(g*g^*)\subseteq(\supp g)(\supp g)^{-1}.
\]
\end{lemma}

\begin{proof}
Continuity and boundedness follow from \cref{prop:l1-banach-algebra}.  For
$x_1,\ldots,x_n\in G$ and $c_1,\ldots,c_n\in\C$, a change of variables gives
\[
 \sum_{i,j}c_i\overline{c_j}(g*g^*)(x_i^{-1}x_j)
 =\int_G\abs{\sum_i c_i g(x_i^{-1}y)}^2\,dy\geq0.
\]
The support assertion follows directly from the convolution integral.
\end{proof}

\section{The Fourier algebra and the spectrum of $\Lp1(G)$}
\label{sec:l1-spectrum}

\subsection{Elementary Fourier analysis}

\begin{proposition}[Fourier transform into $\Czero$]
\label{prop:riemann-lebesgue}
For $f\in\Lp1(G)$, the function $\widehat f$ of
\eqref{eq:fourier-convention} belongs to $\Czero(\dual G)$ and
\[
  \norm{\widehat f}_\infty\leq\norm f_1.
\]
The identities \eqref{eq:fourier-identities} hold.
\end{proposition}

\begin{proof}
The norm bound is immediate.  If $\chi_i\to\chi$ uniformly on compact sets, choose a
compact $K\subseteq G$ with $\int_{K^c}\abs f<\varepsilon$.  Uniform convergence on
$K$ and the bound $\abs{\chi_i-\chi}\leq2$ show
\[
 \abs{\widehat f(\chi_i)-\widehat f(\chi)}
 \leq\int_K\abs f\abs{\chi_i-\chi}+2\varepsilon\longrightarrow2\varepsilon.
\]
Thus $\widehat f$ is continuous.

We prove vanishing at infinity first for $f\in\Cc(G)$.  Let
$\delta=\inf\{\abs{1-z}:z\notin A_{\pi/4}\}>0$.  By
\cref{lem:translation-continuity}, choose a relatively compact identity neighbourhood
$U$ so that $\norm{L_yf-f}_1<\delta\varepsilon$ for $y\in U$.  The polar family
\[
 K_U=\set{\chi:\chi(\overline U)\subseteq\overline{A_{\pi/4}}}
\]
is compact by the Arzel\`a--Ascoli argument in \cref{thm:dual-lca}.  If
$\chi\notin K_U$, choose $y\in\overline U$ with
$\abs{1-\chi(y)}\geq\delta$.  The translation identity yields
\[
 \delta\abs{\widehat f(\chi)}
 \leq\abs{1-\overline{\chi(y)}}\abs{\widehat f(\chi)}
 =\abs{\widehat{L_yf-f}(\chi)}<\delta\varepsilon.
\]
After making the translation estimate on $\overline U$ by shrinking $U$, this proves
$\abs{\widehat f}<\varepsilon$ outside a compact set.  Approximation of an arbitrary
$f\in\Lp1$ by $\Cc(G)$, together with the norm bound, proves the general case.

Finally, Fubini, Haar invariance, and the substitutions $x=ay$ or $x=y^{-1}$ prove the
three identities in \eqref{eq:fourier-identities}.
\end{proof}

\begin{formalization}
The only delicate point in this proof is keeping $\overline U$ inside the region on which
translation is uniformly small.  One convenient Lean statement chooses a compact
neighbourhood $C$ first, proves uniform smallness on some closed $U\subseteq C$, and
defines $K_U$ using that closed $U$.
\end{formalization}

Let $\Delta(\Lp1(G))$ denote the set of nonzero complex-linear multiplicative maps
$\Lambda:\Lp1(G)\to\C$, with the Gelfand topology (pointwise convergence on
$\Lp1(G)$).  Although $\Lp1(G)$ need not be unital, every such map is continuous:
extend it to the unitization by $\widetilde\Lambda(f,c)=\Lambda(f)+c$.  If
$\lambda=\widetilde\Lambda(b)$, then $b-\lambda1$ cannot be invertible, so
$\abs\lambda\leq r(b)\leq\norm b$.  Restriction gives boundedness of $\Lambda$.

\begin{theorem}[Spectrum theorem]
\label{thm:l1-spectrum}
The map
\[
 \nu_G:\dual G\longrightarrow\Delta(\Lp1(G)),\qquad
 \nu_G(\chi)(f)=\widehat f(\chi),
\]
is a homeomorphism.  In particular, the complex characters of the convolution algebra
$\Lp1(G)$ are exactly integration against conjugate group characters.
\end{theorem}

\begin{proof}
The convolution identity makes $\nu_G(\chi)$ multiplicative.  It is nonzero: take a
nonnegative $u\in\Cc(G)$ with nonzero integral and put $f=u\chi$ (using the coercion
$\T\subset\C$); then $\widehat f(\chi)=\int u>0$.  If
$\nu_G(\chi)=\nu_G(\eta)$, then
$\int f(\overline\chi-\overline\eta)=0$ for all $f\in\Lp1$.  Testing against compactly
supported functions shows $\chi=\eta$ almost everywhere, hence everywhere by
continuity.  Thus $\nu_G$ is injective.

Let $\Lambda\in\Delta(\Lp1(G))$.  Choose $f_0$ with $\Lambda(f_0)\neq0$ and define
\[
 q(a)=\frac{\Lambda(L_af_0)}{\Lambda(f_0)}.
\]
The identities
\[
 (L_af)*g=L_a(f*g)=f*(L_ag)
\]
and multiplicativity show, first for $f$ or $g$ on which $\Lambda$ is nonzero and then
for all of $\Lp1$, that
\begin{equation}
  \Lambda(L_af)=q(a)\Lambda(f),\qquad q(ab)=q(a)q(b),\qquad q(1)=1.
  \label{eq:spectrum-q}
\end{equation}
Continuity of translation in $\Lp1$ makes $q$ continuous.  It is bounded because
$\abs{q(a)}\leq\norm\Lambda\norm{f_0}_1/\abs{\Lambda(f_0)}$.  A bounded homomorphism
$G\to\C^\times$ has modulus one: for fixed $a$, boundedness of $q(a)^n$ and
$q(a)^{-n}$ forces $\abs{q(a)}=1$.  Hence $q:G\to\T$ is a character.

Bochner integration of the continuous map $a\mapsto f(a)L_ag$ in $\Lp1$ gives
\[
 f*g=\int_G f(a)L_ag\,da.
\]
Apply $\Lambda$, use \eqref{eq:spectrum-q}, and choose $g$ with
$\Lambda(g)\neq0$:
\[
 \Lambda(f)\Lambda(g)=\Lambda(f*g)
 =\Lambda(g)\int_Gf(a)q(a)\,da.
\]
Thus $\Lambda(f)=\int f q=\widehat f(\overline q)$, proving surjectivity.

It remains to identify the topology.  Let the \emph{transform topology} on $\dual G$ be
the weakest topology making all $\widehat f$, $f\in\Lp1(G)$, continuous.  It is the
topology transported from the Gelfand spectrum.  By
\cref{prop:riemann-lebesgue} it is no finer than the compact--open topology.

We prove joint continuity of $(a,\chi)\mapsto\chi(a)$ when $\dual G$ has the transform
topology.  For fixed $f$, the map
\[
  (a,\chi)\longmapsto\widehat{L_af}(\chi)
\]
is continuous: near $(a_0,\chi_0)$ split the difference into a term bounded by
$\norm{L_af-L_{a_0}f}_1$ and a term controlled by the defining coordinate
$\widehat{L_{a_0}f}$.  At $(a_0,\chi_0)$ choose $f_0$ with
$\widehat f_0(\chi_0)\neq0$.  On a neighbourhood where the denominator stays nonzero,
\[
 \chi(a)=\overline{
   \frac{\widehat{L_af_0}(\chi)}{\widehat f_0(\chi)}}
\]
is therefore continuous.  Given a compact $K\subseteq G$ and an identity neighbourhood
$V\subseteq\T$, joint continuity supplies, for every $a\in K$, a neighbourhood $U_a$ of
$a$ and a transform-open identity neighbourhood $N_a$ such that
$\chi(U_a)\subseteq V$.  A finite subcover of $K$ shows that
$\bigcap_aN_a\subseteq W(K,V)$.  By
\cref{lem:dual-neighbourhood-basis}, every compact--open identity neighbourhood is a
transform-topology neighbourhood.  The two topologies coincide, and $\nu_G$ is a
homeomorphism.
\end{proof}

\subsection{Density and uniqueness}

Put
\[
  \mathcal A(G)=\set{\widehat f:f\in\Lp1(G)}\subseteq\Czero(\dual G).
\]

\begin{proposition}[Density of Fourier transforms]
\label{prop:fourier-dense}
$\mathcal A(G)$ is uniformly dense in $\Czero(\dual G)$.
\end{proposition}

\begin{proof}
It is a subalgebra because $\widehat{f*g}=\widehat f\widehat g$, and it is self-adjoint
because $\widehat{f^*}=\overline{\widehat f}$.  It vanishes nowhere by the construction
used in \cref{thm:l1-spectrum}.  It separates points: if
$\widehat f(\chi)=\widehat f(\eta)$ for every $f$, then the spectrum theorem gives
$\chi=\eta$.  The locally compact Stone--Weierstrass theorem (F6) proves the claim.
\end{proof}

\begin{theorem}[Fourier--Stieltjes uniqueness]
\label{thm:FS-uniqueness}
Let $\sigma$ be a finite complex Radon measure on $\dual G$.  If
\[
  \int_{\dual G}\chi(x)\,d\sigma(\chi)=0
  \qquad\text{for every }x\in G,
\]
then $\sigma=0$.
\end{theorem}

\begin{proof}
For $f\in\Lp1(G)$, Fubini gives
\[
 \int_{\dual G}\widehat f(\chi)\,d\sigma(\chi)
 =\int_G f(x)\left(\int_{\dual G}\overline{\chi(x)}\,d\sigma(\chi)\right)dx=0,
\]
because the inner integral is the assumed transform at $x^{-1}$.  Thus $\sigma$
annihilates $\mathcal A(G)$.  By \cref{prop:fourier-dense} it annihilates all of
$\Czero(\dual G)$, and Riesz--Markov uniqueness gives $\sigma=0$.
\end{proof}

\begin{remark}
\label{rem:noncircular-uniqueness}
This theorem is deliberately asymmetric: measures live on $\dual G$, but the testing
characters are only the evaluations $\chi\mapsto\chi(x)$ for $x\in G$.  Surjectivity of
$G\to\bidual G$ has not been used.  This noncircular form is the decisive input in the
final density argument for Pontryagin duality.
\end{remark}

\section{Positive-type functions and the GNS construction}
\label{sec:positive-type}

\begin{definition}[Positive type]
\label{def:positive-type}
A continuous function $\varphi:G\to\C$ is of \emph{positive type} if, for every
$n\geq1$, $x_1,\ldots,x_n\in G$, and $c_1,\ldots,c_n\in\C$,
\begin{equation}
  \sum_{i,j=1}^n\overline{c_i}c_j\,
     \varphi(x_i^{-1}x_j)\geq0.
  \label{eq:positive-type-finite}
\end{equation}
The inequality means that the displayed complex number is real and nonnegative.
Write $P(G)$ for the convex set of positive-type functions satisfying
$\varphi(1)\leq1$.
\end{definition}

\begin{lemma}[Integral criterion]
\label{lem:positive-integral-criterion}
For a bounded continuous $\varphi$, condition
\eqref{eq:positive-type-finite} is equivalent to
\begin{equation}
  \int_G\int_G \overline{f(x)}f(y)\varphi(x^{-1}y)
       \,d\mu_G(x)d\mu_G(y)\geq0
  \qquad(f\in\Cc(G)).
  \label{eq:positive-type-integral}
\end{equation}
\end{lemma}

\begin{proof}
Assume first the finite condition.  For fixed $f$, all variables lie in the compact set
$K=\supp f$.  Approximate $f$ in $\Lp1$ by compactly supported simple functions and
approximate the continuous kernel $(x,y)\mapsto\varphi(x^{-1}y)$ uniformly on
$K\times K$ by a step kernel obtained from a finite measurable partition.  The resulting
double integral is a limit of expressions
\[
  \sum_{i,j}\overline{d_i}d_j\varphi(t_i^{-1}t_j)
\]
with the cell measures absorbed into $d_i$ and $d_j$; these are nonnegative by
\eqref{eq:positive-type-finite}.  The error is bounded by the uniform kernel error times
$\norm f_1^2$, so the limit is nonnegative.

Conversely, take a normalized bump $h_U$ from \cref{lem:normalized-bump} and apply
\eqref{eq:positive-type-integral} to
$f_U=\sum_i c_iL_{x_i}h_U$.  Expanding the integral gives a finite sum whose $(i,j)$
entry is an average of $\varphi$ over a set shrinking to $x_i^{-1}x_j$.  Continuity of
$\varphi$, compact support of $h_U$, and dominated convergence show that this entry
tends to $\varphi(x_i^{-1}x_j)$.  Passing to the limit gives
\eqref{eq:positive-type-finite}.
\end{proof}

\begin{lemma}[Elementary consequences]
\label{lem:positive-type-elementary}
If $\varphi$ is of positive type, then
\[
 \varphi(1)\in\R_{\geq0},\qquad
 \varphi(x^{-1})=\overline{\varphi(x)},\qquad
 \abs{\varphi(x)}\leq\varphi(1).
\]
In particular, $\varphi(1)=0$ implies $\varphi=0$, and the definition of $P(G)$ is
equivalent to the bound $\norm\varphi_\infty\leq1$.
\end{lemma}

\begin{proof}
The $1\times1$ matrix gives $\varphi(1)\geq0$.  The matrix associated to $1,x$ is
positive semidefinite and therefore Hermitian, giving
$\varphi(x^{-1})=\overline{\varphi(x)}$.  Its determinant is nonnegative:
$\varphi(1)^2-\abs{\varphi(x)}^2\geq0$.  The remaining assertions follow.
\end{proof}

Every character $\chi\in\dual G$ belongs to $P(G)$, because
\[
 \sum_{i,j}\overline{c_i}c_j\chi(x_i^{-1}x_j)
 =\abs{\sum_jc_j\chi(x_j)}^2.
\]
More generally, a positive finite measure $\sigma$ on $\dual G$ has
$\check\sigma$ of positive type by integrating this square.

\subsection{Construction of the cyclic representation}

Let $\varphi$ be of positive type.  On $\Cc(G)$ define
\begin{equation}
 \angles{f,g}_\varphi
 =\int_G\int_G\overline{f(x)}g(y)\varphi(x^{-1}y)\,dx\,dy.
 \label{eq:gns-form}
\end{equation}
This is a positive semidefinite sesquilinear form, conjugate-linear in the first
variable.  Cauchy--Schwarz follows by applying nonnegativity to
$f+zg$ and minimizing over $z\in\C$.  Hence
\[
 N_\varphi=\set{f:\angles{f,f}_\varphi=0}
\]
is a linear subspace and the quotient $\Cc(G)/N_\varphi$ is an inner-product space.
Let $H_\varphi$ be its Hilbert completion.

\begin{proposition}[GNS representation]
\label{prop:gns}
There is a strongly continuous unitary representation
\[
 U^\varphi:G\longrightarrow\mathcal U(H_\varphi)
\]
and a cyclic vector $\xi_\varphi\in H_\varphi$ such that
\begin{equation}
  \varphi(a)=\angles{\xi_\varphi,U^\varphi_a\xi_\varphi}
  \qquad(a\in G).
  \label{eq:gns-coefficient}
\end{equation}
The linear span of $\{U^\varphi_a\xi_\varphi:a\in G\}$ is dense in $H_\varphi$.
\end{proposition}

\begin{proof}
Define $U_a[f]=[L_af]$.  Haar invariance and a change of variables in
\eqref{eq:gns-form} give
$\angles{L_af,L_ag}_\varphi=\angles{f,g}_\varphi$, so this is well-defined and extends
to a unitary operator on $H_\varphi$.  The translation identity gives
$U_aU_b=U_{ab}$.  Also
\begin{equation}
  \norm{[f]}_\varphi^2
  \leq\norm\varphi_\infty\norm f_1^2.
  \label{eq:gns-l1-bound}
\end{equation}
Thus \cref{lem:translation-continuity} implies
$\norm{U_a[f]-[f]}_\varphi\to0$ for the dense set of classes of $\Cc$ functions.
Unitarity extends this to every vector, proving strong continuity.

Define the conjugate-linear functional on $\Cc(G)$
\[
  \Phi(f)=\int_G\overline{f(x)}\varphi(x^{-1})\,dx.
\]
If $f\in N_\varphi$, Cauchy--Schwarz gives
$\angles{f,h}_\varphi=0$ for all $h\in\Cc(G)$.  Put
\[
 F_f(y)=\int_G\overline{f(x)}\varphi(x^{-1}y)\,dx.
\]
This is continuous and $\angles{f,h}_\varphi=\int h(y)F_f(y)\,dy$.  Testing against all
$h$ gives $F_f=0$, and in particular $\Phi(f)=F_f(1)=0$.  Thus $\Phi$ descends to the
quotient.

For the approximate identity $h_U$,
\[
 \Phi(f)=\lim_U\angles{f,h_U}_\varphi.
\]
Indeed, the right side averages the continuous function $F_f$ over a set shrinking to
$1$.  Moreover, \eqref{eq:gns-l1-bound} and $\norm{h_U}_1=1$ give
$\norm{[h_U]}_\varphi\leq\sqrt{\norm\varphi_\infty}$.  Therefore
\[
 \abs{\Phi(f)}\leq\sqrt{\norm\varphi_\infty}\norm{[f]}_\varphi.
\]
The functional extends continuously to $H_\varphi$.  By the Hilbert-space Riesz theorem
there is a unique $\xi_\varphi$ such that
$\Phi(f)=\angles{[f],\xi_\varphi}_\varphi$.

Replacing $f$ by a translate, and using unitarity, gives for every $f\in\Cc(G)$
\begin{equation}
 \angles{[f],U_a\xi_\varphi}_\varphi
 =\int_G\overline{f(y)}\varphi(y^{-1}a)\,dy.
 \label{eq:gns-translate-vector}
\end{equation}
Indeed, this follows from
$\angles{[f],U_a\xi}=\angles{U_{a^{-1}}[f],\xi}$ and the change of variables $y=ax$.
Integrating \eqref{eq:gns-translate-vector} against $h(a)\,da$ shows, by comparison with
\eqref{eq:gns-form}, that the Bochner integral
\[
  \int_G h(a)U_a\xi_\varphi\,da
\]
equals the vector $[h]$.  Therefore a vector orthogonal to every $U_a\xi_\varphi$ is
orthogonal to every class $[h]$, hence is zero.  This proves cyclicity.  Finally,
\[
 \angles{\xi_\varphi,[h]}
 =\overline{\Phi(h)}
 =\int_Gh(a)\varphi(a)\,da,
\]
where \cref{lem:positive-type-elementary} was used in the last equality.  Comparing this
with $[h]=\int h(a)U_a\xi_\varphi\,da$ shows that
$a\mapsto\angles{\xi_\varphi,U_a\xi_\varphi}$ equals $\varphi(a)$ almost everywhere.
The left side is continuous by strong continuity and the right side is continuous by
hypothesis, so full support of Haar measure upgrades the equality to every $a$.
\end{proof}

\begin{remark}[Measurable positive-type functions]
\label{rem:positive-continuous-rep}
The same proof begins with an essentially bounded measurable function satisfying the
integral positivity condition.  It produces the continuous representative
$a\mapsto\angles{\xi,U_a\xi}$.  This fact is needed below to make the positive unit ball
weak-$*$ closed without building continuity into a non-closed predicate.
\end{remark}

\subsection{Extreme points}

\begin{lemma}[Extreme coefficient criterion]
\label{lem:extreme-irreducible}
Let $\varphi\in P(G)$ with $\varphi(1)=1$.  Then $\varphi$ is an extreme point of
$P(G)$ if and only if its cyclic GNS representation is irreducible.
\end{lemma}

\begin{proof}
Suppose first that $W\subseteq H_\varphi$ is a closed invariant subspace and let $P$ be
its orthogonal projection.  Since the representation is unitary, $W^\perp$ is invariant
and $P$ commutes with every $U_a$.  Put
\begin{align*}
 \varphi_1(a)&=\angles{P\xi,U_aP\xi},&
 \varphi_2(a)&=\angles{(1-P)\xi,U_a(1-P)\xi}.
\end{align*}
Both are of positive type and $\varphi=\varphi_1+\varphi_2$.  If $\varphi$ is extreme,
the elementary extreme-point criterion for the cone slice
$\{\psi:\psi(1)\leq1\}$ says that
$\varphi_1=\lambda\varphi$ for some $0\leq\lambda\leq1$.  Polarizing the equality of
matrix coefficients after replacing $a$ by $b^{-1}a$ gives
\[
 \angles{U_b\xi,PU_a\xi}=\lambda\angles{U_b\xi,U_a\xi}.
\]
Cyclicity implies $P=\lambda I$.  Since $P$ is a projection, $\lambda\in\{0,1\}$;
thus $W$ is trivial.

Conversely, if the representation is reducible, take a nontrivial invariant $W$.  By
cyclicity neither $P\xi$ nor $(1-P)\xi$ is zero.  After normalizing $\varphi_1$ and
$\varphi_2$ by their positive values at $1$, the equality
$\varphi=\varphi_1+\varphi_2$ becomes a nontrivial convex decomposition in $P(G)$.
Hence $\varphi$ is not extreme.
\end{proof}

\begin{lemma}[Abelian Schur lemma]
\label{lem:abelian-schur}
Every irreducible strongly continuous unitary representation of an abelian group on a
nonzero complex Hilbert space is one-dimensional.  It has the form
$U_a=\chi(a)I$ for a unique $\chi\in\dual G$.
\end{lemma}

\begin{proof}
For fixed $a$, the unitary $U_a$ commutes with every $U_b$.  Every spectral projection
of the normal operator $U_a$ therefore commutes with the representation, so its range is
invariant.  Irreducibility forces every spectral projection to be $0$ or $I$; the spectral
measure is concentrated at one point, and $U_a=\chi(a)I$.  The representation law makes
$\chi$ multiplicative, unitarity gives $\abs\chi=1$, and strong continuity gives
continuity of $\chi$.  Since now every closed subspace is invariant, irreducibility forces
the Hilbert space to have dimension one.
\end{proof}

\begin{theorem}[Extreme points of $P(G)$]
\label{thm:extreme-characters}
The extreme points of $P(G)$ are the zero function and the characters in $\dual G$.
\end{theorem}

\begin{proof}
The zero function is extreme because a positive-type function vanishing at $1$ is zero.
If a nonzero $\varphi\in P(G)$ is extreme, then necessarily $\varphi(1)=1$; otherwise
\[
 \varphi=\varphi(1)\frac{\varphi}{\varphi(1)}+(1-\varphi(1))0
\]
is a nontrivial convex decomposition.  By
\cref{lem:extreme-irreducible,lem:abelian-schur}, its GNS representation is
one-dimensional, say $U_a=\chi(a)I$.  Equation \eqref{eq:gns-coefficient} and
$\norm\xi^2=\varphi(1)=1$ give $\varphi=\chi$.

Conversely, the GNS space of a character is one-dimensional: in
\eqref{eq:gns-form},
\[
 \norm{[f]}_\chi^2
 =\abs{\int_Gf(x)\chi(x)\,dx}^2.
\]
It is irreducible, so the character is extreme by
\cref{lem:extreme-irreducible}.
\end{proof}

\section{Bochner's theorem}
\label{sec:bochner}

We prove Bochner's theorem rather than use it as a disguised form of duality.  Regard an
essentially bounded positive-type function as an element of
$\Lp\infty(G)=\Lp1(G)^*$ with its weak-$*$ topology.

\begin{lemma}[Weak-$*$ compact positive ball]
\label{lem:P-weakstar-compact}
The set $P(G)$, identified with its $\Lp\infty$ classes, is weak-$*$ compact and convex.
\end{lemma}

\begin{proof}
For $f\in\Cc(G)$ put
\[
 k_f(t)=\int_G\overline{f(x)}f(xt)\,dx.
\]
Then $k_f\in\Lp1(G)$ and, by Fubini,
\[
 \int_G\int_G\overline{f(x)}f(y)\varphi(x^{-1}y)\,dx\,dy
 =\int_G\varphi(t)k_f(t)\,dt.
\]
Thus integral positivity is the intersection, over $f\in\Cc(G)$, of weak-$*$ closed
real half-spaces in the closed unit ball of $\Lp\infty$.  Banach--Alaoglu makes that
intersection compact.  By \cref{rem:positive-continuous-rep}, every class in it has a
unique continuous positive-type representative; uniqueness uses full support.  Hence it
is precisely $P(G)$.  Convexity is immediate.
\end{proof}

\begin{theorem}[Bochner]
\label{thm:bochner}
A continuous function $\varphi:G\to\C$ is of positive type if and only if there is a
unique finite positive Radon measure $\sigma$ on $\dual G$ such that
\begin{equation}
  \varphi(x)=\int_{\dual G}\chi(x)\,d\sigma(\chi).
  \label{eq:bochner}
\end{equation}
Moreover $\sigma(\dual G)=\varphi(1)$.
\end{theorem}

\begin{proof}
After dividing by $\varphi(1)$, it suffices to treat $P(G)$.  Let
\[
 \mathcal M=\set{\check\sigma:
   \sigma\geq0,\ \sigma(\dual G)\leq1}\subseteq\Lp\infty(G).
\]
The positive unit ball of finite measures on $\dual G=\Czero(\dual G)^*$ is weak-$*$
compact.  The map $\sigma\mapsto\check\sigma$ is continuous from that weak-$*$ topology
to the weak-$*$ topology of $\Lp\infty(G)$.  Indeed, for $h\in\Lp1(G)$, Fubini writes
\[
 \int_Gh(x)\check\sigma(x)\,dx
 =\int_{\dual G}\left(\int_Gh(x)\chi(x)\,dx\right)d\sigma(\chi),
\]
and the inner function of $\chi$ belongs to $\Czero(\dual G)$ by
\cref{prop:riemann-lebesgue}.  Consequently $\mathcal M$ is compact, hence weak-$*$
closed.  It is convex and contained in $P(G)$, since an integral of characters is of
positive type.

The set $\mathcal M$ contains $0$ and every character (use the Dirac measures).
By \cref{thm:extreme-characters}, it contains every extreme point of $P(G)$.
Krein--Milman says that $P(G)$ is the closed convex hull of its extreme points, so
$P(G)\subseteq\mathcal M$.  Thus equality holds.  Uniqueness of $\sigma$ is
\cref{thm:FS-uniqueness}, and evaluating \eqref{eq:bochner} at $1$ gives its total mass.
\end{proof}

\begin{formalization}
The weak-$*$ closure proof should be carried out on $\Lp\infty$ equivalence classes; only
after compactness is established should one choose the continuous GNS representative.
This avoids attempting to prove that continuity is weak-$*$ closed.  The use of the
spectral theorem in \cref{lem:abelian-schur} is generic Hilbert-space functional analysis
and may be packaged as a separate Schur lemma.
\end{formalization}

\section{Construction of the dual Haar measure and Fourier inversion}
\label{sec:inversion}

Let
\[
 V(G)=\Span_\C P(G),\qquad V^1(G)=V(G)\cap\Lp1(G).
\]
By linearity of Bochner's theorem, for every $f\in V(G)$ there is a finite complex
measure $\sigma_f$ on $\dual G$ such that
\begin{equation}
  f(x)=\int_{\dual G}\chi(x)\,d\sigma_f(\chi).
  \label{eq:sigma-f}
\end{equation}
To see that this measure is independent of a presentation of $f$ as a linear combination
of positive-type functions, subtract two candidate measures and apply
\cref{thm:FS-uniqueness}.  Thus $f\mapsto\sigma_f$ is a well-defined linear map.  When
$f$ is of positive type, $\sigma_f$ is positive.

\subsection{A symmetric measure identity}

\begin{lemma}
\label{lem:symmetric-measure-identity}
For $f,g\in V^1(G)$,
\begin{equation}
  \widehat g\,\sigma_f=\widehat f\,\sigma_g
  \quad\text{as finite complex measures on }\dual G.
  \label{eq:symmetric-measure-identity}
\end{equation}
\end{lemma}

\begin{proof}
Both products are finite measures because Fourier transforms are bounded.  It suffices by
\cref{thm:FS-uniqueness} to compare their inverse Fourier--Stieltjes transforms.  Fubini
and \eqref{eq:sigma-f} give
\begin{align*}
 \int_{\dual G}\chi(x)\widehat g(\chi)\,d\sigma_f(\chi)
 &=\int_Gg(y)
    \left(\int_{\dual G}\chi(xy^{-1})\,d\sigma_f(\chi)\right)dy\\
 &=\int_Gg(y)f(xy^{-1})\,dy=(g*f)(x).
\end{align*}
The other measure has transform $(f*g)(x)$, equal to this because $G$ is abelian.
\end{proof}

Choose the approximate identity $h_U$ of \cref{prop:approximate-identity} and put
\begin{equation}
  e_U=h_U*h_U^*.
  \label{eq:positive-approx-identity}
\end{equation}
By \cref{lem:convolution-positive-type}, $e_U\in V^1(G)$ is of positive type, so
$\sigma_{e_U}$ is positive.  Moreover
\begin{equation}
 \widehat{e_U}=\abs{\widehat{h_U}}^2\longrightarrow1
 \quad\text{uniformly on compact subsets of }\dual G.
 \label{eq:positive-approx-transform}
\end{equation}

\subsection{The auxiliary multiplier class}

Define $\mathcal F$ to be the set of $u\in\Cb(\dual G)$ for which there exists a finite
complex measure $\tau_u$ on $\dual G$ satisfying
\begin{equation}
  u\,\sigma_f=\widehat f\,\tau_u
  \qquad\text{for every }f\in V^1(G).
  \label{eq:F-defining}
\end{equation}
Products of a function and a measure mean multiplication of the measure by that bounded
Borel function.

\begin{lemma}[Multiplier calculus]
\label{lem:multiplier-calculus}
The class $\mathcal F$ has the following properties.
\begin{enumerate}[label=(\roman*)]
\item The measure $\tau_u$ is unique.  If $u\geq0$, then $\tau_u\geq0$.
\item $\mathcal F$ is a $\Cb(\dual G)$-module, and
\[
 \tau_{u+v}=\tau_u+\tau_v,\quad
 \tau_{cu}=c\tau_u,\quad
 \tau_{au}=a\tau_u
\]
whenever the expressions make sense.
\item For $\eta\in\dual G$, set
$(T_\eta u)(\chi)=u(\eta^{-1}\chi)$ and let $T_\eta\tau$ be pushforward under
$\chi\mapsto\eta\chi$.  Then
$T_\eta u\in\mathcal F$ and $\tau_{T_\eta u}=T_\eta\tau_u$.
\item If $g\in V^1(G)$, then $\widehat g\in\mathcal F$ and
$\tau_{\widehat g}=\sigma_g$.
\item $\Cc(\dual G)\subseteq\mathcal F$.
\end{enumerate}
\end{lemma}

\begin{proof}
Apply \eqref{eq:F-defining} with $f=e_U$.  For $a\in\Cc(\dual G)$,
\begin{equation}
 \int a\,d\tau_u
 =\lim_U\int a\widehat{e_U}\,d\tau_u
 =\lim_U\int au\,d\sigma_{e_U}.
 \label{eq:tau-recovery}
\end{equation}
The first limit follows from uniform convergence on $\supp a$.  The last expression is
determined by $u$, proving uniqueness.  If $u\geq0$ and $a\geq0$, it is nonnegative
because $\sigma_{e_U}$ is positive; Riesz--Markov then gives $\tau_u\geq0$.

Linearity follows immediately from \eqref{eq:F-defining} and uniqueness.  If
$a\in\Cb(\dual G)$, multiply both sides of \eqref{eq:F-defining} by $a$ to obtain
$(au)\sigma_f=\widehat f(a\tau_u)$; this proves the module assertion.

For translation, define the modulation $(M_\eta f)(x)=\eta(x)f(x)$.  Directly from
\eqref{eq:sigma-f} and \eqref{eq:fourier-convention},
\begin{equation}
 \sigma_{M_\eta f}=T_\eta\sigma_f,
 \qquad \widehat{M_\eta f}=T_\eta\widehat f.
 \label{eq:modulation-translation}
\end{equation}
Apply \eqref{eq:F-defining} to $M_{\eta^{-1}}f$ and push the resulting equality forward
by $T_\eta$.  This gives
$(T_\eta u)\sigma_f=\widehat f(T_\eta\tau_u)$.

Assertion (iv) is exactly \cref{lem:symmetric-measure-identity}: for all $f$,
$\widehat g\,\sigma_f=\widehat f\,\sigma_g$.

Finally let $\gamma\in\Cc(\dual G)$ and $K=\supp\gamma$.  By
\eqref{eq:positive-approx-transform}, choose $U$ with
$\abs{\widehat{e_U}-1}<1/2$ on $K$.  Hence $\widehat{e_U}$ is nonzero on an open
neighbourhood of $K$.  The function
\[
 a(\chi)=
 \begin{cases}
   \gamma(\chi)/\widehat{e_U}(\chi),&\widehat{e_U}(\chi)\neq0,\\
   0,&\widehat{e_U}(\chi)=0
 \end{cases}
\]
is continuous and compactly supported: near the zero set of $\widehat{e_U}$ it is zero
because that set is disjoint from $K=\supp\gamma$.  Thus $a\in\Cb(\dual G)$ and
$\gamma=a\widehat{e_U}$.  Assertions (ii) and (iv) imply $\gamma\in\mathcal F$.
\end{proof}

\subsection{The dual measure}

For $\gamma\in\Cc(\dual G)$ define
\begin{equation}
  \Lambda(\gamma)=\tau_\gamma(\dual G).
  \label{eq:Lambda-dual-measure}
\end{equation}
This is well-defined by \cref{lem:multiplier-calculus}.

\begin{proposition}[Dual Haar measure]
\label{prop:dual-haar-measure}
$\Lambda$ is a nonzero positive translation-invariant linear functional on
$\Cc(\dual G)$.  Therefore there is a nonzero Haar measure $\mu_{\dual G}$ such that
\begin{equation}
  \Lambda(\gamma)=\int_{\dual G}\gamma(\chi)\,d\mu_{\dual G}(\chi)
  \qquad(\gamma\in\Cc(\dual G)).
  \label{eq:dual-haar-riesz}
\end{equation}
\end{proposition}

\begin{proof}
Linearity and positivity follow from parts (i)--(ii) of
\cref{lem:multiplier-calculus}.  A positive functional on $\Cc$ is locally bounded:
if $K$ is compact, choose $b\in\Cc$, $b\geq1$ on $K$; decomposing a complex
$\gamma$ supported in $K$ into four nonnegative parts bounds
$\abs{\Lambda(\gamma)}$ by a fixed multiple of
$\norm\gamma_\infty\Lambda(b)$.  Riesz--Markov therefore supplies a positive Radon
measure satisfying \eqref{eq:dual-haar-riesz}.

It is nonzero.  Choose nonzero $g\in\Cc(G)$ and put $f=g*g^*$.  Then $f$ is nonzero,
integrable, and of positive type, so $\sigma_f$ is a nonzero positive measure.  Choose
$a\in\Cc(\dual G)$, $a\geq0$, with $\int a\,d\sigma_f>0$.  Since
$\widehat f\in\mathcal F$ and $a\widehat f\in\Cc(\dual G)$,
\[
 \Lambda(a\widehat f)=\tau_{a\widehat f}(\dual G)
 =\int a\,d\sigma_f>0.
\]

For $\eta\in\dual G$, part (iii) of \cref{lem:multiplier-calculus} gives
$\Lambda(T_\eta\gamma)=(T_\eta\tau_\gamma)(\dual G)=\Lambda(\gamma)$.  Thus the
representing measure is translation invariant and nonzero, hence is Haar.
\end{proof}

\begin{theorem}[Fourier inversion on $V^1$]
\label{thm:fourier-inversion-V}
For every $f\in V^1(G)$,
\begin{equation}
  \sigma_f=\widehat f\,\mu_{\dual G},
  \qquad \widehat f\in\Lp1(\dual G),
  \qquad
  f(x)=\int_{\dual G}\widehat f(\chi)\chi(x)\,d\mu_{\dual G}(\chi).
  \label{eq:inversion-V}
\end{equation}
The Haar measure $\mu_{\dual G}$ is uniquely determined by this property.
\end{theorem}

\begin{proof}
We first prove, for every $u\in\mathcal F$, that
\begin{equation}
  \tau_u=u\,\mu_{\dual G}.
  \label{eq:tau-density}
\end{equation}
If $a\in\Cc(\dual G)$, then $au\in\Cc(\dual G)$ and the module identity gives
\[
 \int a(\chi)u(\chi)\,d\mu_{\dual G}(\chi)
 =\Lambda(au)=\tau_{au}(\dual G)=\int a\,d\tau_u.
\]
Riesz--Markov uniqueness proves \eqref{eq:tau-density}.  Apply it to
$u=\widehat f$, for which $\tau_u=\sigma_f$ by
\cref{lem:multiplier-calculus}.  Since $\sigma_f$ is finite,
$\widehat f\in\Lp1(\dual G)$, and substituting its density in
\eqref{eq:sigma-f} proves inversion.

If $\nu$ is another Haar measure with the same inversion property, then
$\widehat f\,\nu=\sigma_f=\widehat f\,\mu_{\dual G}$ for every $f\in V^1$.  Given
$a\in\Cc(\dual G)$, choose $e_U$ so that $\widehat{e_U}$ has no zero on
$\supp a$.  Testing the measure equality for $e_U$ against
$a/\widehat{e_U}$ gives $\int a\,d\nu=\int a\,d\mu_{\dual G}$.  Hence $\nu$ equals
$\mu_{\dual G}$.
\end{proof}

\begin{corollary}[Positivity and scaling]
\label{cor:dual-measure-scaling}
If $f\in V^1(G)$ is of positive type, then $\widehat f\geq0$ everywhere.  If the original
Haar measure is replaced by $c\mu_G$ with $c>0$, the dual Haar measure selected by
\cref{thm:fourier-inversion-V} is replaced by $c^{-1}\mu_{\dual G}$.
\end{corollary}

\begin{proof}
For positive-type $f$, the measure $\sigma_f=\widehat f\mu_{\dual G}$ is positive, so
$\widehat f\geq0$ almost everywhere.  It is continuous and Haar measure has full
support, so the inequality is pointwise.  Under scaling by $c$, the Fourier transform
is multiplied by $c$; multiplying the dual measure by $c^{-1}$ leaves the inversion
integral unchanged.  Uniqueness gives the assertion.
\end{proof}

\section{Plancherel theory and localized Fourier transforms}
\label{sec:plancherel}

The only purpose of the full $\Lp2$ theorem in the biduality proof is to manufacture a
nonzero Fourier transform supported in a prescribed nonempty open set.  We nevertheless
record the complete extension because it is a central part of the theory and fixes the
normalization of dual Haar measure.

\begin{theorem}[Parseval on the dense domain]
\label{thm:parseval-dense}
If $f\in\Lp1(G)\cap\Lp2(G)$, then $\widehat f\in\Lp2(\dual G)$ and
\begin{equation}
  \int_G\abs{f(x)}^2\,d\mu_G(x)
  =\int_{\dual G}\abs{\widehat f(\chi)}^2\,d\mu_{\dual G}(\chi).
  \label{eq:parseval-dense}
\end{equation}
\end{theorem}

\begin{proof}
Set $k=f*f^*$.  By \cref{prop:l1-banach-algebra}, $k\in\Lp1(G)$ is continuous; by
\cref{lem:convolution-positive-type}, it is of positive type.  Hence $k\in V^1(G)$.
Furthermore
\[
  k(1)=\norm f_2^2,
  \qquad \widehat k=\widehat f\,\widehat{f^*}=\abs{\widehat f}^2.
\]
Apply inversion \eqref{eq:inversion-V} to $k$ at $1$.
\end{proof}

\begin{lemma}[Density of the Fourier range]
\label{lem:plancherel-range-dense}
The image of $\Lp1(G)\cap\Lp2(G)$ under the Fourier transform is dense in
$\Lp2(\dual G)$.
\end{lemma}

\begin{proof}
Suppose $H\in\Lp2(\dual G)$ is orthogonal to the image.  Fix
$f\in\Lp1(G)\cap\Lp2(G)$.  For every $x\in G$, the function
$\chi\mapsto\chi(x)\widehat f(\chi)$ is the Fourier transform of
$L_{x^{-1}}f$ and is again in the image.  Therefore
\[
  \int_{\dual G}\overline{H(\chi)}\widehat f(\chi)\chi(x)
       \,d\mu_{\dual G}(\chi)=0
  \qquad(x\in G).
\]
The density $\overline H\widehat f$ belongs to $\Lp1(\dual G)$ by Cauchy--Schwarz, so
it defines a finite measure.  Fourier--Stieltjes uniqueness
\cref{thm:FS-uniqueness} gives
\begin{equation}
  \overline H\widehat f=0\quad\text{almost everywhere}
  \qquad\text{for every }f\in\Lp1\cap\Lp2.
  \label{eq:H-times-transform-zero}
\end{equation}

If $H$ were nonzero, regularity of the finite measure
$\abs H^2\mu_{\dual G}$ would give a compact set $K$ on which
$\int_K\abs H^2>0$.  Choose $e_U$ from
\eqref{eq:positive-approx-identity} so that
$\abs{\widehat{e_U}-1}<1/2$ on $K$.  Since
$e_U\in\Cc(G)\subseteq\Lp1\cap\Lp2$, equation
\eqref{eq:H-times-transform-zero} applied to $e_U$ contradicts the choice of $K$.
Thus $H=0$, and the orthogonal complement of the image is trivial.
\end{proof}

\begin{theorem}[Plancherel]
\label{thm:plancherel}
The Fourier transform extends uniquely to a unitary operator
\[
  \mathcal P_G:\Lp2(G,\mu_G)\longrightarrow
     \Lp2(\dual G,\mu_{\dual G}).
\]
For $f,g\in\Lp2(G)$,
\begin{equation}
  \int_G\overline{f(x)}g(x)\,dx
  =\int_{\dual G}\overline{\mathcal P_Gf(\chi)}
      \mathcal P_Gg(\chi)\,d\chi.
  \label{eq:plancherel-inner-product}
\end{equation}
\end{theorem}

\begin{proof}
The space $\Lp1\cap\Lp2$ is dense in $\Lp2$, for instance because $\Cc(G)$ is dense.
By \cref{thm:parseval-dense}, Fourier transform is an isometry there, so it extends
uniquely to an isometry on $\Lp2$.  Its range is closed and is dense by
\cref{lem:plancherel-range-dense}, hence is all of $\Lp2(\dual G)$.  Polarization gives
\eqref{eq:plancherel-inner-product}.
\end{proof}

\begin{lemma}[$\Lp2$ product--convolution formula]
\label{lem:l2-product-convolution}
If $a,b\in\Lp2(G)$, then $ab\in\Lp1(G)$ and
\begin{equation}
  \widehat{ab}=(\mathcal P_Ga)*(\mathcal P_Gb),
  \label{eq:l2-product-convolution}
\end{equation}
where the convolution on the right is the continuous bounded representative of the
convolution of two $\Lp2(\dual G)$ functions.
\end{lemma}

\begin{proof}
H\"older gives $ab\in\Lp1$.  Fix $\eta\in\dual G$ and put
$d(x)=\overline{b(x)}\eta(x)$.  Then $d\in\Lp2(G)$ and a direct calculation on the dense
domain $\Lp1\cap\Lp2$, followed by $\Lp2$ approximation, gives
\[
  (\mathcal P_Gb)(\chi^{-1}\eta)
  =\overline{(\mathcal P_Gd)(\chi)}.
\]
Consequently Parseval yields
\begin{align*}
 ((\mathcal P_Ga)*(\mathcal P_Gb))(\eta)
 &=\int_{\dual G}(\mathcal P_Ga)(\chi)
      (\mathcal P_Gb)(\chi^{-1}\eta)\,d\chi\\
 &=\int_{\dual G}(\mathcal P_Ga)(\chi)
      \overline{(\mathcal P_Gd)(\chi)}\,d\chi\\
 &=\int_Ga(x)\overline{d(x)}\,dx
   =\int_Ga(x)b(x)\overline{\eta(x)}\,dx
   =\widehat{ab}(\eta).
\end{align*}
Both sides have canonical continuous representatives, so the pointwise identity holds
everywhere.
\end{proof}

\begin{proposition}[Fourier transform localized in an open set]
\label{prop:localized-fourier}
Let $U$ be a nonempty open subset of $\dual G$.  There exists
$F\in\Lp1(G)$ such that $\widehat F$ is nonzero and
\[
  \supp(\widehat F)\subseteq U.
\]
\end{proposition}

\begin{proof}
Haar measure is positive on $U$ and inner regular, so choose compact
$K\subseteq U$ with $0<\mu_{\dual G}(K)<\infty$.  Compactness of $K$ and continuity of
multiplication give an identity neighbourhood $W$ with $WK\subseteq U$.  Choose a
compact neighbourhood $B$ of $1$ contained in $W$; then
$0<\mu_{\dual G}(B)<\infty$.  The indicator functions $\one_B,\one_K$ belong to
$\Lp2(\dual G)$.  Their convolution
\[
  q=\one_B*\one_K
\]
has a continuous representative, is supported in the compact set $BK\subseteq U$, and
is nonzero because
\[
  \int_{\dual G}q
  =\mu_{\dual G}(B)\mu_{\dual G}(K)>0.
\]
By Plancherel choose $a,b\in\Lp2(G)$ whose transforms are $\one_B$ and $\one_K$.
Then $F=ab\in\Lp1(G)$ and
\cref{lem:l2-product-convolution} gives $\widehat F=q$.
\end{proof}

\begin{formalization}
For indicator functions it is safest to choose compact measurable $B,K$ with positive
finite measure and prove that the continuous convolution representative vanishes off
$BK$.  The theorem is used once with the LCA group $\dual G$ and an open subset of
$\bidual G$.
\end{formalization}

\section{Pontryagin biduality}
\label{sec:biduality}

We now prove that the continuous homomorphism $\alpha_G$ of
\cref{prop:evaluation-map} is a topological isomorphism.  The proof is divided into four
properties so that each can become an independent Lean theorem.

\subsection{Characters separate points}

\begin{proposition}[Injectivity]
\label{prop:alpha-injective}
$\alpha_G:G\to\bidual G$ is injective.
\end{proposition}

\begin{proof}
Suppose $z\in G$ satisfies $\chi(z)=1$ for every $\chi\in\dual G$.  If $z\neq1$, choose
an identity neighbourhood $U$ with $zU\cap U=\varnothing$.  Choose a relatively compact
identity neighbourhood $V$ with $VV^{-1}\subseteq U$, and a nonzero
$g\in\Cc(G)$ supported in $V$.  The function
\[
  f=\frac{g*g^*}{(g*g^*)(1)}
\]
is nonzero, belongs to $V^1(G)$, is of positive type, satisfies $f(1)=1$, and has support
in $U$.

Fourier inversion gives
\[
 f(z^{-1}x)=\int_{\dual G}\widehat f(\chi)
       \chi(z)^{-1}\chi(x)\,d\chi=f(x)
\]
for every $x$.  Thus $L_zf=f$.  But $\supp(L_zf)=z\supp f\subseteq zU$ is disjoint from
$\supp f\subseteq U$, impossible for a nonzero function.  Hence $z=1$.
\end{proof}

\begin{corollary}
\label{cor:characters-separate}
For $x\neq y$ in $G$, there is $\chi\in\dual G$ with $\chi(x)\neq\chi(y)$.
\end{corollary}

\begin{proof}
Apply \cref{prop:alpha-injective} to $x^{-1}y\neq1$.
\end{proof}

\subsection{The original topology is recovered from characters}

For compact $K\subseteq\dual G$ and an identity neighbourhood $V\subseteq\T$, put
\[
  K^{\triangleleft}_V
  =\set{x\in G:\chi(x)\in V\text{ for every }\chi\in K}
  =\alpha_G^{-1}(W(K,V)).
\]

\begin{proposition}[Topological embedding]
\label{prop:alpha-embedding}
The sets $K^{\triangleleft}_V$ form a neighbourhood basis at $1$ in $G$.  Consequently
$\alpha_G$ is a homeomorphism from $G$ onto its image in $\bidual G$.
\end{proposition}

\begin{proof}
Each $K^{\triangleleft}_V$ is a neighbourhood of $1$ because $\alpha_G$ is continuous.
We prove cofinality in the other direction.  Let $U$ be an identity neighbourhood.
As in \cref{prop:alpha-injective}, choose a continuous positive-type
$f\in V^1(G)$ with
\[
  f(1)=1,\qquad \supp f\subseteq U.
\]
By \cref{cor:dual-measure-scaling,thm:fourier-inversion-V},
$\widehat f\geq0$ and
\[
  1=f(1)=\int_{\dual G}\widehat f(\chi)\,d\chi.
\]
Choose $0<\varepsilon<1/3$.  Regularity of the finite measure
$\widehat f\,\mu_{\dual G}$ gives a compact $K\subseteq\dual G$ such that
\begin{equation}
  \int_{K^c}\widehat f(\chi)\,d\chi<\varepsilon.
  \label{eq:embedding-tail}
\end{equation}
Let $V=\set{z\in\T:\abs{z-1}<\varepsilon}$.  If
$x\in K^{\triangleleft}_V$, inversion and \eqref{eq:embedding-tail} give
\begin{align*}
 \abs{f(x)-1}
 &\leq\int_K\widehat f(\chi)\abs{\chi(x)-1}\,d\chi
   +\int_{K^c}\widehat f(\chi)\abs{\chi(x)-1}\,d\chi\\
 &<\varepsilon\int_K\widehat f+2\varepsilon<3\varepsilon<1.
\end{align*}
Thus $f(x)\neq0$, so $x\in\supp f\subseteq U$.  We have found
$K^{\triangleleft}_V\subseteq U$, proving the basis assertion.  Since the subspace
neighbourhoods of $1$ in $\alpha_G(G)$ are precisely the intersections with
$W(K,V)$, the inverse of $\alpha_G$ on its image is continuous.
\end{proof}

\subsection{The image is closed}

\begin{lemma}[Dense locally compact subspaces]
\label{lem:dense-locally-compact-open}
If $X$ is a dense locally compact subspace of a Hausdorff space $Y$, then $X$ is open in
$Y$.
\end{lemma}

\begin{proof}
Fix $x\in X$.  Choose a compact neighbourhood $K$ of $x$ in $X$ and a relatively open
set $V\subseteq X$ with $x\in V\subseteq K$.  Write $V=X\cap O$ for an open
$O\subseteq Y$.  The set $K$ is compact in $Y$, hence closed.  If $O\setminus K$ were
nonempty, it would be a nonempty open subset of $Y$ and would meet the dense set $X$.
But
\[
  X\cap(O\setminus K)=V\setminus K=\varnothing,
\]
a contradiction.  Hence $x\in O\subseteq K\subseteq X$, so $X$ is open.
\end{proof}

\begin{proposition}[Closed image]
\label{prop:alpha-closed}
The subgroup $\alpha_G(G)$ is closed in $\bidual G$.
\end{proposition}

\begin{proof}
Let $H=\overline{\alpha_G(G)}$.  By
\cref{prop:alpha-embedding}, $\alpha_G(G)$ is locally compact in its subspace topology;
it is dense in the Hausdorff space $H$.  By
\cref{lem:dense-locally-compact-open}, it is open in $H$.  An open subgroup is closed,
since its complement is the union of the other open cosets.  Thus $\alpha_G(G)$ is both
dense and closed in $H$, so it equals $H$ and is closed in $\bidual G$.
\end{proof}

\subsection{The image is all of the bidual}

\begin{proposition}[Surjectivity]
\label{prop:alpha-surjective}
$\alpha_G:G\to\bidual G$ is surjective.
\end{proposition}

\begin{proof}
Suppose not.  By \cref{prop:alpha-closed},
\[
  \Omega=\bidual G\setminus\alpha_G(G)
\]
is a nonempty open set.  Apply \cref{prop:localized-fourier} to the LCA group
$\dual G$.  There is $\Phi\in\Lp1(\dual G)$ whose Fourier transform on $\bidual G$ is
nonzero and supported in $\Omega$.  It therefore vanishes at every evaluation character:
\begin{equation}
  0=\widehat\Phi(\alpha_G(x))
  =\int_{\dual G}\Phi(\chi)\overline{\alpha_G(x)(\chi)}\,d\chi
  =\int_{\dual G}\Phi(\chi)\chi(x^{-1})\,d\chi.
  \label{eq:surjectivity-zero-transform}
\end{equation}
The finite complex measure $\Phi\,\mu_{\dual G}$ therefore has zero inverse
Fourier--Stieltjes transform at every $x\in G$.  The asymmetric uniqueness theorem
\cref{thm:FS-uniqueness}---which does not assume biduality---implies
$\Phi\,\mu_{\dual G}=0$.  Hence $\Phi=0$ almost everywhere and
$\widehat\Phi=0$, contradicting its construction.  Thus $\Omega$ is empty.
\end{proof}

\begin{theorem}[Pontryagin duality]
\label{thm:pontryagin}
For every LCA group $G$, evaluation is a natural isomorphism of topological groups
\begin{equation}
  \alpha_G:G\xrightarrow{\ \sim\ }\bidual G,
  \qquad \alpha_G(x)(\chi)=\chi(x).
  \label{eq:pontryagin-isomorphism}
\end{equation}
It is natural in $G$, and the triangle identity
$(\dual{\alpha_G})\alpha_{\dual G}=\id_{\dual G}$ holds.
\end{theorem}

\begin{proof}
Continuity is \cref{prop:evaluation-map}; injectivity, continuity of the inverse, and
surjectivity are
\cref{prop:alpha-injective,prop:alpha-embedding,prop:alpha-surjective}.  Naturality and
the triangle identity are \cref{prop:naturality}.
\end{proof}

\begin{corollary}[Mutual normalization of Haar measures]
\label{cor:mutual-haar-normalization}
Start with $\mu_G$, construct $\mu_{\dual G}$ by
\cref{prop:dual-haar-measure}, and then construct the dual of $\mu_{\dual G}$.  Under
$\alpha_G:G\simeq\bidual G$, the resulting Haar measure on $\bidual G$ is exactly
$\mu_G$, not merely a scalar multiple.
\end{corollary}

\begin{proof}
Transport the second dual measure along $\alpha_G$; it is $c\mu_G$ for some $c>0$ by
Haar uniqueness.  For $F\in\Lp1(\dual G)\cap\Lp2(\dual G)$, the adjoint formula for the
unitary Plancherel operator gives
\[
  \widehat F(\alpha_G(x))=(\mathcal P_G^{-1}F)(x^{-1}).
\]
Apply Plancherel for the group $\dual G$ and then change variables by inversion:
\[
 \norm F_{2,\mu_{\dual G}}^2
 =\norm{\widehat F}_{2,\mu_{\bidual G}}^2
 =c\norm{\mathcal P_G^{-1}F}_{2,\mu_G}^2
 =c\norm F_{2,\mu_{\dual G}}^2.
\]
Choose $F\neq0$ to obtain $c=1$.
\end{proof}

\section{Annihilators, subgroups, quotients, and exactness}
\label{sec:annihilators}

For $S\subseteq G$ and $T\subseteq\dual G$ define
\begin{align*}
 S^\perp&=\set{\chi\in\dual G:\chi(s)=1\text{ for all }s\in S},\\
 T^\perp&=\set{x\in G:\chi(x)=1\text{ for all }\chi\in T}.
\end{align*}
The second definition uses the evaluation embedding, now known to identify $G$ with its
bidual.

\begin{lemma}[Elementary annihilator laws]
\label{lem:annihilator-elementary}
For arbitrary subsets $S,S_i\subseteq G$:
\begin{enumerate}[label=(\alph*)]
\item $S^\perp$ is a closed subgroup of $\dual G$;
\item inclusion reverses: $S\subseteq T$ implies $T^\perp\subseteq S^\perp$;
\item $S^\perp=(\overline{\angles S})^\perp$, where $\overline{\angles S}$ is the
closed subgroup generated by $S$;
\item $(\bigcup_iS_i)^\perp=\bigcap_iS_i^\perp$.
\end{enumerate}
The analogous statements hold for subsets of $\dual G$.
\end{lemma}

\begin{proof}
The annihilator is the intersection over $s\in S$ of the kernels of the continuous
evaluation homomorphisms $\chi\mapsto\chi(s)$, hence is a closed subgroup.  The other
claims follow from the definitions, continuity of characters, and the fact that a
character trivial on $S$ is trivial on the subgroup generated by $S$ and then on its
closure.
\end{proof}

Let $H$ be a closed subgroup of $G$, and let $q:G\to G/H$ be the quotient map.  The
quotient is again an LCA group.

\begin{theorem}[Dual of a quotient]
\label{thm:dual-quotient}
Precomposition with $q$ is a topological isomorphism
\begin{equation}
  \dual q:\dual{(G/H)}\xrightarrow{\ \sim\ }H^\perp.
  \label{eq:dual-quotient}
\end{equation}
\end{theorem}

\begin{proof}
A character $\psi$ of $G/H$ pulls back to a character of $G$ trivial on $H$.  Conversely,
if $\chi\in H^\perp$, then $\bar\chi(gH)=\chi(g)$ is well-defined and is a homomorphism.
It is continuous because $q$ is a quotient map and $\bar\chi\circ q=\chi$.  These
constructions are inverse, proving algebraic bijectivity.  Continuity of $\dual q$ is
functoriality of the compact--open topology.

For continuity of the inverse we use the following compact lifting fact.  Every compact
$C\subseteq G/H$ is contained in $q(K)$ for some compact $K\subseteq G$.  Indeed, choose
a compact identity neighbourhood $B\subseteq G$.  The open sets
$yq(B^\circ)$ cover $C$; choose finitely many, lift their translating elements to
$x_1,\ldots,x_n\in G$, and take $K=\bigcup_ix_iB$.

Now a compact--open neighbourhood $W(C,V)$ in $\dual{(G/H)}$ contains the inverse image
of $W(K,V)\cap H^\perp$ under \eqref{eq:dual-quotient}, because $C\subseteq q(K)$.
Thus the inverse is continuous at the identity and hence everywhere.
\end{proof}

\begin{theorem}[Double annihilator]
\label{thm:double-annihilator}
If $H$ is a closed subgroup of $G$, then
\begin{equation}
  (H^\perp)^\perp=H.
  \label{eq:double-annihilator}
\end{equation}
For an arbitrary subset $S\subseteq G$,
\begin{equation}
  (S^\perp)^\perp=\overline{\angles S}.
  \label{eq:double-annihilator-subset}
\end{equation}
\end{theorem}

\begin{proof}
The inclusion $H\subseteq(H^\perp)^\perp$ is immediate.  If $x\notin H$, then
$q(x)\neq1$ in the LCA group $G/H$.  By
\cref{cor:characters-separate}, choose $\psi\in\dual{(G/H)}$ with
$\psi(q(x))\neq1$.  Its pullback $\psi\circ q$ belongs to $H^\perp$ and does not vanish
at $x$.  Thus $x\notin(H^\perp)^\perp$, proving \eqref{eq:double-annihilator}.  The
general formula follows from \cref{lem:annihilator-elementary} by replacing $S$ with
$\overline{\angles S}$.
\end{proof}

\begin{corollary}[Lattice formulas]
\label{cor:annihilator-lattice}
For closed subgroups $H,K\leq G$,
\begin{align}
 (\overline{HK})^\perp&=H^\perp\cap K^\perp,
 \label{eq:annihilator-join}\\
 (H\cap K)^\perp&=\overline{H^\perp K^\perp}.
 \label{eq:annihilator-meet}
\end{align}
The same formulas hold for arbitrary families, with closed generated subgroup replacing
the finite product.
\end{corollary}

\begin{proof}
The first identity is immediate from the definition.  Let
$A=\overline{H^\perp K^\perp}$.  Applying the first identity in the dual group and then
double annihilation gives
\[
 A^\perp=(H^\perp)^\perp\cap(K^\perp)^\perp=H\cap K.
\]
Taking annihilators once more gives \eqref{eq:annihilator-meet}.  The family version is
identical.
\end{proof}

\subsection{Extension of characters}

We first isolate a useful formal consequence of biduality.

\begin{lemma}[The dual reflects isomorphisms]
\label{lem:dual-reflects-iso}
If $u:A\to B$ is a continuous homomorphism of LCA groups and
$\dual u:\dual B\to\dual A$ is a topological isomorphism, then $u$ is a topological
isomorphism.
\end{lemma}

\begin{proof}
Naturality \eqref{eq:naturality} rearranges to
\begin{equation}
  u=\alpha_B^{-1}\circ\bidual u\circ\alpha_A.
  \label{eq:dual-reflects-iso}
\end{equation}
If $\dual u$ is an isomorphism, so is its dual $\bidual u$; the two evaluation maps are
isomorphisms by \cref{thm:pontryagin}.  Hence the right side is an isomorphism.
\end{proof}

\begin{theorem}[Character extension and dual of a subgroup]
\label{thm:character-extension}
Let $i:H\hookrightarrow G$ be the inclusion of a closed subgroup.  Restriction
\[
  \dual i:\dual G\longrightarrow\dual H,
  \qquad\chi\longmapsto\chi|_H
\]
is a continuous open surjection with kernel $H^\perp$.  Equivalently, every continuous
character of $H$ extends to a continuous character of $G$, and restriction induces a
topological isomorphism
\begin{equation}
  \dual G/H^\perp\xrightarrow{\ \sim\ }\dual H.
  \label{eq:dual-subgroup}
\end{equation}
\end{theorem}

\begin{proof}
The kernel is visibly $H^\perp$.  Put $N=H^\perp$, $Q=\dual G/N$, and let
$\bar r:Q\to\dual H$ be the injective map induced by restriction.  It is continuous by
the quotient property.  We prove that its dual is an isomorphism.

Apply \cref{thm:dual-quotient} to the closed subgroup $N\leq\dual G$.  It identifies
$\dual Q$ with $N^\perp\subseteq\bidual G$.  By
\cref{thm:double-annihilator},
\[
  N^\perp=(H^\perp)^\perp=\alpha_G(H).
\]
Consequently the map
\[
 \theta:H\longrightarrow\dual Q,
 \qquad \theta(h)(\chi N)=\chi(h)
\]
is a topological isomorphism.  For $h\in H$ and $\chi N\in Q$,
\begin{align*}
 ((\dual{\bar r})(\alpha_H(h)))(\chi N)
 &=\alpha_H(h)(\bar r(\chi N))\\
 &=(\chi|_H)(h)=\chi(h)=\theta(h)(\chi N).
\end{align*}
Thus $(\dual{\bar r})\circ\alpha_H=\theta$ is an isomorphism.  Since $\alpha_H$ is an
isomorphism, $\dual{\bar r}$ is one too.  By
\cref{lem:dual-reflects-iso}, $\bar r$ is a topological isomorphism.  In particular
restriction is surjective and, as the composite of the open quotient map
$\dual G\to Q$ with an isomorphism, is open.
\end{proof}

\begin{theorem}[Exactness]
\label{thm:exactness}
For a closed subgroup $H\leq G$, the strict exact sequence
\[
  1\longrightarrow H\xrightarrow{i}G\xrightarrow{q}G/H\longrightarrow1
\]
dualizes to the strict exact sequence
\begin{equation}
  1\longrightarrow\dual{(G/H)}
  \xrightarrow{\dual q}\dual G
  \xrightarrow{\dual i}\dual H\longrightarrow1.
  \label{eq:dual-exact-sequence}
\end{equation}
Here $\dual q$ is a closed embedding and $\dual i$ is an open quotient map.
\end{theorem}

\begin{proof}
By \cref{thm:dual-quotient}, $\dual q$ identifies the first group with $H^\perp$.
By \cref{thm:character-extension}, this subgroup is the kernel of the surjective open map
$\dual i$.  These statements give algebraic exactness and the asserted topology.
\end{proof}

\begin{proposition}[Kernel and closure of the dual image]
\label{prop:dual-kernel-image}
For any continuous homomorphism $f:G\to H$,
\begin{align}
 \ker(\dual f)&=(\overline{f(G)})^\perp,
 \label{eq:kernel-dual-map}\\
 \overline{(\dual f)(\dual H)}&=(\ker f)^\perp.
 \label{eq:closure-image-dual-map}
\end{align}
If $f$ occurs as the inclusion or quotient map in a strict exact sequence, the closure in
\eqref{eq:closure-image-dual-map} may be omitted.
\end{proposition}

\begin{proof}
A character $\psi\in\dual H$ lies in the kernel precisely when it is trivial on $f(G)$,
and continuity makes this equivalent to triviality on its closure.  For the second
formula, $x\in G$ annihilates the image of $\dual f$ exactly when
\[
  \psi(f(x))=1\quad\text{for every }\psi\in\dual H.
\]
Characters separate points of $H$, so this is equivalent to $f(x)=1$.  Thus the
annihilator of the image is $\ker f$; take annihilators again and use
\cref{thm:double-annihilator}.  The last assertion follows from
\cref{thm:dual-quotient,thm:character-extension}.
\end{proof}

\section{Categorical and structural consequences}
\label{sec:categorical}

\begin{theorem}[Contravariant equivalence]
\label{thm:anti-equivalence}
Pontryagin duality is a contravariant equivalence on the category of LCA groups and
continuous homomorphisms.  Its object map is $G\mapsto\dual G$, its morphism map is
$f\mapsto\dual f$, and its unit is the natural isomorphism $\alpha:\id\Rightarrow\dual{(
\dual{\mathord{-}})}$.
\end{theorem}

\begin{proof}
Contravariant identity and composition laws were proved in
\cref{sec:evaluation}.  The maps $\alpha_G$ are isomorphisms by
\cref{thm:pontryagin}, are natural by \eqref{eq:naturality}, and obey the triangle
identity \eqref{eq:triangle}.  These data are precisely a dual equivalence.

Explicitly, the functor is faithful: if $\dual f=\dual g$, then for every $x\in G$ and
$\psi\in\dual H$,
$\psi(f(x))=(\dual f\,\psi)(x)=(\dual g\,\psi)(x)=\psi(g(x))$; characters separate
points, so $f=g$.  It is full in the contravariant sense.  Given a continuous
$u:\dual H\to\dual G$, define
\begin{equation}
 f=\alpha_H^{-1}\circ\dual u\circ\alpha_G:G\longrightarrow H.
 \label{eq:recover-map-from-dual}
\end{equation}
For $x\in G$ and $\psi\in\dual H$,
\[
 \psi(f(x))=(\dual u(\alpha_G(x)))(\psi)
 =\alpha_G(x)(u(\psi))=u(\psi)(x),
\]
so $\dual f=u$ by extensionality.
\end{proof}

\begin{theorem}[Compact/discrete equivalence]
\label{thm:compact-discrete-iff}
For an LCA group $G$,
\begin{align*}
 G\text{ is compact}&\quad\Longleftrightarrow\quad\dual G\text{ is discrete},\\
 G\text{ is discrete}&\quad\Longleftrightarrow\quad\dual G\text{ is compact}.
\end{align*}
Thus Pontryagin duality restricts to a dual equivalence between compact abelian groups
and discrete abelian groups.
\end{theorem}

\begin{proof}
The forward implications are \cref{prop:compact-discrete-basic}.  If $\dual G$ is
discrete, then $\bidual G$ is compact, and $G\simeq\bidual G$ is compact.  If
$\dual G$ is compact, then $\bidual G$ is discrete, and hence so is $G$.
\end{proof}

\begin{corollary}
An LCA group is both compact and discrete if and only if it is finite.  In that case its
dual is finite of the same cardinality.
\end{corollary}

\begin{proof}
A compact discrete space has a finite cover by singleton open sets.  Conversely a finite
Hausdorff group is compact and discrete.  For finite $G$, biduality gives
$\abs G=\abs{\bidual G}$; finite character theory, or the cyclic decomposition in
\cref{ex:finite}, gives $\abs{\dual G}=\abs G$.
\end{proof}

\begin{proposition}[Closed subgroup tests]
\label{prop:annihilator-compact-open}
For a closed subgroup $H\leq G$:
\begin{enumerate}[label=(\alph*)]
\item if $H$ is open, then $H^\perp$ is compact;
\item if $H$ is compact, then $H^\perp$ is open in $\dual G$;
\item $G/H$ is compact if and only if $H^\perp$ is discrete;
\item $G/H$ is discrete if and only if $H^\perp$ is compact.
\end{enumerate}
\end{proposition}

\begin{proof}
If $H$ is open, $G/H$ is discrete, so
$H^\perp\simeq\dual{(G/H)}$ is compact.  If $H$ is compact, $\dual H$ is discrete and
restriction $\dual G\to\dual H$ is continuous; the kernel $H^\perp$, the inverse image
of the open singleton $\{1\}$, is open.  The last two statements follow from
\cref{thm:dual-quotient,thm:compact-discrete-iff}.
\end{proof}

\subsection{Products and sums}

\begin{proposition}[Finite products]
\label{prop:dual-finite-product}
For LCA groups $G$ and $H$, restriction to the two factors is a topological isomorphism
\[
  \dual{(G\times H)}\xrightarrow{\ \sim\ }\dual G\times\dual H.
\]
Its inverse sends $(\chi,\psi)$ to $(g,h)\mapsto\chi(g)\psi(h)$.
\end{proposition}

\begin{proof}
Every $(g,h)$ factors as $(g,1)(1,h)$, which proves that the displayed maps are
algebraic inverses.  They are continuous by compact--open functoriality.  For continuity
of the inverse, if $K\subseteq G\times H$ is compact, its projections $K_G,K_H$ are
compact.  Choose identity neighbourhoods $V_G,V_H\subseteq\T$ with
$V_GV_H\subseteq V$.  Then
\[
  W(K_G,V_G)\times W(K_H,V_H)
\]
maps into $W(K,V)$.  This proves continuity at the identity.
\end{proof}

\begin{proposition}[Discrete sums and compact products]
\label{prop:dual-sum-product}
Let $(D_i)_{i\in I}$ be discrete abelian groups and $(K_i)_{i\in I}$ compact abelian
groups.  Then
\begin{align}
 \dual{\left(\bigoplus_{i\in I}D_i\right)}
   &\simeq\prod_{i\in I}\dual{D_i},
 \label{eq:dual-direct-sum}\\
 \dual{\left(\prod_{i\in I}K_i\right)}
   &\simeq\bigoplus_{i\in I}\dual{K_i}.
 \label{eq:dual-compact-product}
\end{align}
The first isomorphism is between compact groups, and the second between discrete groups.
\end{proposition}

\begin{proof}
A homomorphism from the algebraic direct sum is uniquely a family of homomorphisms from
the factors, and continuity is automatic because the sum is discrete.  Compact subsets
of a discrete space are finite, so the compact--open topology is pointwise convergence
on finitely many elements; this is exactly the product topology.  This proves
\eqref{eq:dual-direct-sum}.

For the second formula, a finite-support family of characters defines a character of the
product.  Conversely, let $\chi:\prod_iK_i\to\T$ be continuous.  The inverse image of
$A_{\pi/2}$ contains a basic identity neighbourhood
$\prod_iU_i$, with $U_i=K_i$ outside a finite set $F$.  For $j\notin F$, the entire
coordinate subgroup $K_j$ maps into $A_{\pi/2}$; by
\cref{lem:nss-circle}, it maps trivially.  Hence $\chi$ depends on only the coordinates
in $F$.  This proves algebraic bijectivity.  Both sides are discrete: the product is
compact, so its dual is discrete, and every $\dual{K_i}$ is discrete and the direct sum
is given the discrete topology.  Thus the bijection is a homeomorphism.
\end{proof}

\subsection{A general Fourier inversion criterion}

\begin{theorem}
\label{thm:general-fourier-inversion}
Let $f\in\Lp1(G)$.  If $\widehat f\in\Lp1(\dual G)$, then
\begin{equation}
  f(x)=\int_{\dual G}\widehat f(\chi)\chi(x)\,d\chi
  \label{eq:general-inversion}
\end{equation}
for almost every $x$.  The right side is a continuous function vanishing at infinity;
hence \eqref{eq:general-inversion} holds at every point at which $f$ has that continuous
representative.  In particular, the Fourier transform is injective on $\Lp1(G)$.
\end{theorem}

\begin{proof}
Let
$g(x)=\int\widehat f(\chi)\chi(x)\,d\chi$.  Dominated convergence and the
Riemann--Lebesgue argument on the dual group show $g\in\Czero(G)$.  With $e_U$ from
\eqref{eq:positive-approx-identity}, Fubini and the already proved inversion formula for
$e_U$ give
\[
 (f*e_U)(x)
 =\int_{\dual G}\widehat f(\chi)\widehat{e_U}(\chi)\chi(x)\,d\chi.
\]
The left side converges to $f$ in $\Lp1$.  Since
$\abs{\widehat{e_U}}\leq1$ and $\widehat{e_U}\to1$ pointwise, dominated convergence
shows that the right side converges to $g$ uniformly in $x$.  On each compact
$K\subseteq G$, uniform convergence implies convergence in $\Lp1(K)$, while global
$\Lp1$ convergence also implies convergence in $\Lp1(K)$.  Uniqueness of the local
$\Lp1$ limit gives $f=g$ almost everywhere on every compact set; Radon inner regularity
then gives equality almost everywhere on $G$.  If $\widehat f=0$, then $g=0$, proving
injectivity.
\end{proof}

\section{Canonical examples}
\label{sec:examples}

In the following examples additive notation is used inside formulas, while the abstract
duality functor remains multiplicative.

\begin{example}[$\Z$ and the circle]
\label{ex:Z-circle}
The map
\[
  \T\longrightarrow\dual\Z,
  \qquad z\longmapsto(n\mapsto z^n)
\]
is a topological isomorphism.  Indeed, a homomorphism from $\Z$ is uniquely determined
by the image of $1$, and every such homomorphism is continuous because $\Z$ is discrete.
The displayed map is continuous; it is a continuous bijection from compact $\T$ to the
Hausdorff space $\dual\Z$, hence a homeomorphism.  Dualizing and applying
\cref{thm:pontryagin} gives the canonical isomorphism
\[
  \dual\T\simeq\Z,
  \qquad n\longmapsto(z\mapsto z^n).
\]
The topology on $\dual\T$ is discrete by compactness of $\T$.
\end{example}

\begin{example}[The real line]
\label{ex:R}
For $\xi\in\R$ put
\[
  \chi_\xi(x)=\exp(2\pi i\xi x).
\]
Then $\xi\mapsto\chi_\xi$ is a topological isomorphism
$\R\simeq\dual\R$.

To prove surjectivity without invoking biduality, let $\chi:\R\to\T$ be a continuous
additive character.  On a sufficiently small interval $(-\delta,\delta)$, the image lies
in an arc on which a continuous argument is unique.  Write
$\chi(x)=e^{2\pi ia(x)}$ there, with $a(0)=0$.  Whenever
$x,y,x+y\in(-\delta,\delta)$, the difference
$a(x+y)-a(x)-a(y)$ is an integer and varies continuously; at $(0,0)$ it is zero, so it
is zero throughout the connected small domain.  Thus $a$ is locally additive.  For
arbitrary $x$, choose $n$ with $x/n\in(-\delta,\delta)$ and set
$\widetilde a(x)=na(x/n)$.  Local additivity shows that this is independent of $n$ and
is a continuous additive lift of $\chi$.  Every continuous additive
$\widetilde a:\R\to\R$ has the form $\widetilde a(x)=\xi x$: first for rational $x$,
then for real $x$ by density.  Hence $\chi=\chi_\xi$.

The parameter map is continuous for the compact--open topology by uniform continuity of
$(\xi,x)\mapsto e^{2\pi i\xi x}$ on compact rectangles.  Conversely, uniform closeness
of $\chi_\xi$ to $1$ on $[-1,1]$ first forces $\abs\xi<1/2$ (otherwise some
$x\in[-1,1]$ maps to $-1$), after which the principal argument at $x=1$ controls
$\abs\xi$.  Thus the inverse parameter map is continuous.
\end{example}

\begin{example}[Euclidean groups]
\label{ex:Rn}
By \cref{prop:dual-finite-product,ex:R},
\[
 \dual{\R^n}\simeq\R^n,
 \qquad
 \xi\longmapsto\bigl(x\mapsto e^{2\pi i\langle\xi,x\rangle}\bigr).
\]
Under Lebesgue measure on both sides this is the normalization used in
\eqref{eq:fourier-convention}; it is self-dual, so no $(2\pi)^n$ factor occurs in
inversion.
\end{example}

\begin{example}[Finite cyclic groups]
\label{ex:finite-cyclic}
For $C_n=\Z/n\Z$ with the discrete topology,
\[
 C_n\xrightarrow{\ \sim\ }\dual{C_n},
 \qquad k\longmapsto
 \left(m\longmapsto\exp\left(\frac{2\pi i km}{n}\right)\right).
\]
This is an algebraic bijection between finite discrete groups and hence a topological
isomorphism.  Counting measure on $C_n$ is dual to $(1/n)$ times counting measure on its
dual with our inversion convention.
\end{example}

\begin{example}[Finite abelian groups]
\label{ex:finite}
The structure theorem writes a finite abelian group noncanonically as a finite product of
cyclic groups.  The preceding example and finite-product duality give a noncanonical
isomorphism $A\simeq\dual A$.  In contrast, the evaluation isomorphism
$A\simeq\bidual A$ is canonical and natural.  This distinction should be preserved in a
formal API: a finite group generally has no preferred isomorphism with its first dual.
\end{example}

\begin{example}[Tori and lattices]
For the $n$-torus $\T^n$ and the lattice $\Z^n$,
\[
  \dual{\T^n}\simeq\Z^n,
  \qquad \dual{\Z^n}\simeq\T^n.
\]
More generally, if $\Lambda\leq\R^n$ is a lattice, quotient duality gives
\[
  \dual{(\R^n/\Lambda)}\simeq
  \Lambda^\perp
  =\set{\xi\in\R^n:\langle\xi,\lambda\rangle\in\Z
      \text{ for every }\lambda\in\Lambda},
\]
the usual dual lattice, with the discrete topology.
\end{example}

\begin{example}[Infinite sums and products]
For any set $I$, \cref{prop:dual-sum-product} specializes to
\[
  \dual{\left(\bigoplus_I\Z\right)}\simeq\T^I,
  \qquad
  \dual{(\T^I)}\simeq\bigoplus_I\Z.
\]
The first source is discrete and the first target compact; the second source is compact
even for uncountable $I$, and every continuous character depends on finitely many
coordinates.
\end{example}

\section{A Mathlib-oriented implementation plan}
\label{sec:mathlib-plan}

This section translates the proof into declarations and dependency layers.  It is not a
claim that the proposed declaration names already exist.  The items marked ``existing''
were checked against \mathlib~v4.32.0; all others are proposed additions.

\subsection{Existing kernel}

\begin{longtable}{>{\raggedright\arraybackslash}p{0.40\textwidth}
                  >{\raggedright\arraybackslash}p{0.52\textwidth}}
\toprule
Existing declaration or module & Role \\
\midrule
\endhead
\path{PontryaginDual} & Definition as \texttt{A ->\_t* Circle}, with induced
compact--open topology. \\
\path{PontryaginDual.ext}, \path{one_apply} & Extensionality and evaluation simp
normal forms. \\
\path{PontryaginDual.map}, \path{map_apply}, \path{map_comp} & Contravariant
precomposition. \\
\path{PontryaginDual.mapHom} & Continuity of the map operation in a locally compact
codomain parameter. \\
\path{ContinuousMonoidHom.isInducing_toContinuousMap} & Transfers compact--open
topology results from continuous maps. \\
\path{ContinuousMonoidHom.instContinuousEval} & Joint evaluation for a locally
compact domain. \\
\path{ContinuousMonoidHom.continuous_of_continuous_uncurry} & Currying joint
evaluation to construct $\alpha_G$. \\
\path{ContinuousMonoidHom.locallyCompactSpace_of_hasBasis} & Local compactness of
the character group from a countable arc basis in $\T$. \\
\path{Circle.centeredArc}, \path{Circle.mem_centeredArc},
\path{Circle.hasBasis_centeredArc_div_two_pow} & The small-arc basis. \\
\path{Circle.eq_one_of_forall_pow_mem_centeredArc_pi_div_two} & No small
subgroups for the circle. \\
\path{Mathlib.MeasureTheory.Measure.Haar.Basic} & Haar measure and invariance. \\
\path{Mathlib.MeasureTheory.Group.IntegralConvolution} and
\path{Mathlib.Analysis.Convolution} & Convolution infrastructure. \\
\path{Mathlib.Analysis.Convex.KreinMilman} & Krein--Milman. \\
\path{Mathlib.Analysis.Normed.Module.WeakDual} & Weak dual topology and
Banach--Alaoglu infrastructure. \\
\path{Mathlib.Topology.ContinuousMap.StoneWeierstrass} & Density of the Fourier
algebra. \\
\path{Mathlib.MeasureTheory.Integral.RieszMarkovKakutani.Basic} & Representation of
positive functionals by Radon measures. \\
\bottomrule
\end{longtable}

The source already gives the dual three central instances:
\texttt{T2Space}, \texttt{IsTopologicalGroup}, and
\texttt{LocallyCompactSpace}.  It also gives compactness for a discrete source and
discreteness for a compact source.  These facts should be reused, not reproved inside the
biduality file.

\subsection{Recommended file and declaration layers}

\begin{longtable}{>{\raggedright\arraybackslash}p{0.08\textwidth}
                  >{\raggedright\arraybackslash}p{0.31\textwidth}
                  >{\raggedright\arraybackslash}p{0.51\textwidth}}
\toprule
Layer & Suggested module & Principal new declarations \\
\midrule
\endhead
L0 & \path{Topology/Algebra/PontryaginDual/Evaluation} &
\texttt{eval}, \texttt{eval\_apply}, group-hom laws, continuity, naturality, triangle
identity. \\
L1 & \path{Analysis/Fourier/LCA/ApproximateIdentity} & Normalized $\Cc$ bumps,
translation continuity in $\Lp p$, approximate identities, compact-uniform convergence
of their character integrals. \\
L2 & \path{Analysis/Fourier/LCA/Spectrum} & LCA convolution algebra, Riemann--Lebesgue,
$\Lp1$ spectrum equivalence, transform topology equals compact--open, density in
$\Czero$, Fourier--Stieltjes uniqueness. \\
L3 & \path{Analysis/Fourier/LCA/PositiveType} & Finite-matrix positive type, integral
criterion, elementary bounds, convolution-square examples. \\
L4 & \path{Analysis/Fourier/LCA/GNS} & Semidefinite quotient, Hilbert completion,
strongly continuous cyclic representation, extreme/irreducible criterion, abelian Schur
lemma. \\
L5 & \path{Analysis/Fourier/LCA/Bochner} & Weak-$*$ compact positive ball,
Krein--Milman proof, representing positive measure, uniqueness. \\
L6 & \path{Analysis/Fourier/LCA/Inversion} & $V^1(G)$, $\sigma_f$, multiplier class
$\mathcal F$, dual Haar measure, inversion and scaling. \\
L7 & \path{Analysis/Fourier/LCA/Plancherel} & Dense-domain Parseval, unitary extension,
$\Lp2$ product--convolution, localized Fourier transform. \\
L8 & \path{Topology/Algebra/PontryaginDual/Duality} & Injectivity, neighbourhood-basis
recovery, closed image, surjectivity, \texttt{ContinuousMulEquiv} biduality. \\
L9 & \path{Topology/Algebra/PontryaginDual/Exact} & Annihilators, quotient dual,
double annihilator, character extension, strict exactness. \\
L10 & \path{Topology/Algebra/PontryaginDual/Examples} & $\Z$, $\T$, $\R$, finite
cyclic groups, finite products, discrete sums/compact products. \\
\bottomrule
\end{longtable}

This layering keeps the existing lightweight definition file from importing measure
theory and functional analysis.  In particular, the definition and functorial map remain
usable without importing the proof of biduality.

\subsection{Evaluation map skeleton}

The first implementation milestone is independent of the analytic proof.  A typical
constructor skeleton is:

\begin{lstlisting}
namespace PontryaginDual

noncomputable def eval (G : Type*)
    [CommGroup G] [TopologicalSpace G]
    [IsTopologicalGroup G] [LocallyCompactSpace G] :
    G ->_t* PontryaginDual (PontryaginDual G) where
  toFun x :=
    { toFun := fun chi => chi x
      map_one' := by simp
      map_mul' := by intro chi psi; rfl
      continuous_toFun := continuous_eval_const x }
  map_one' := by ext chi; exact chi.map_one
  map_mul' x y := by ext chi; exact chi.map_mul x y
  continuous_toFun := by
    apply ContinuousMonoidHom.continuous_of_continuous_uncurry
    -- discharge with joint evaluation, after swapping factors
    exact ...

@[simp] theorem eval_apply (x : G) (chi : PontryaginDual G) :
    eval G x chi = chi x := rfl

end PontryaginDual
\end{lstlisting}

The actual names of the evaluation lemmas should be resolved with \texttt{\#check} in the
target branch.  It is preferable to prove continuity through the existing
\texttt{ContinuousEval} instance rather than unfold the compact--open topology.

\subsection{Suggested theorem inventory}

The following inventory is fine-grained enough for autoformalization.  Each line should
be a declaration with only earlier lines as substantive dependencies.

\begin{enumerate}[label=\textbf{P\arabic*.},leftmargin=4em]
\item Normalized $\Cc$ bump subordinate to an identity neighbourhood.
\item $\Lp p$ translation isometry and continuity.
\item Approximate identity convergence; uniform character convergence on compact sets.
\item $\Lp1$ convolution norm, associativity, involution, and convolution square of
positive type.
\item Fourier norm bound, continuity, translation/involution/convolution identities.
\item Riemann--Lebesgue vanishing at infinity.
\item Every nonzero multiplicative functional on the unitization of $\Lp1$ is bounded.
\item Spectrum classification $\Delta(\Lp1(G))\simeq\dual G$.
\item Joint evaluation for the transform topology and equality with compact--open.
\item Stone--Weierstrass density of $\{\widehat f\}$ in $\Czero(\dual G)$.
\item Fourier--Stieltjes uniqueness in the asymmetric form of
\cref{thm:FS-uniqueness}.
\item Equivalence of finite-matrix and integral positive type.
\item GNS construction and continuous cyclic coefficient representative.
\item Extreme coefficient if and only if irreducible GNS representation.
\item Abelian Schur lemma and classification of extreme points.
\item Weak-$*$ compactness of $P(G)$.
\item Bochner existence and uniqueness.
\item Symmetric measure identity \eqref{eq:symmetric-measure-identity}.
\item Positive approximate net $e_U$ and recovery formula
\eqref{eq:tau-recovery}.
\item Multiplier class closure, translation, and inclusion of $\Cc(\dual G)$.
\item Construction, nontriviality, and invariance of dual Haar measure.
\item Inversion on $V^1$ and uniqueness of dual normalization.
\item Dense Parseval isometry and density of its range.
\item Plancherel and localized-transform lemma.
\item Evaluation injective.
\item Polar sets recover the original neighbourhood filter.
\item Evaluation image closed.
\item Evaluation image surjective.
\item Assemble \texttt{ContinuousMulEquiv}; prove naturality simp lemmas.
\item Quotient duality, double annihilator, character extension, and exactness.
\end{enumerate}

\subsection{Representation choices and known pressure points}

\paragraph{Multiplicative versus additive APIs.}
The public theorem should first be proved for the multiplicative
\texttt{PontryaginDual}.  Fourier integration uses the coercion
\texttt{Circle -> Complex}.  Only after the multiplicative development is stable should
additive wrappers be generated or proved.  Premature \texttt{to\_additive} use is likely
to create brittle names around convolution and character evaluation.

\paragraph{Positive type.}
Use the finite-matrix predicate as the public definition: it is algebraic, does not
mention Haar measure, and applies to arbitrary topological groups.  Add the integral
criterion as a theorem under LCA and Haar hypotheses.  For matrices indexed by
\texttt{Fin n}, the expression can be written as a quadratic form; this makes the
$2\times2$ bound and closure under limits reusable.

\paragraph{Complex measures.}
Positive measures can use \texttt{Measure}; linear combinations in $V(G)$ require
complex vector measures (the \texttt{MeasureTheory.VectorMeasure} hierarchy) or an
explicit four-positive-measure representation.  The latter is initially easier for
Bochner but makes uniqueness and multiplier identities verbose.  A thin local structure
containing a complex vector measure together with finite-variation evidence is likely the
cleanest interface.

\paragraph{Weak-$*$ topology.}
State positivity on $\Lp\infty$ classes by the family of inequalities
$\operatorname{re}\langle\varphi,k_f\rangle\geq0$ and
$\operatorname{im}\langle\varphi,k_f\rangle=0$.  These are manifestly weak-$*$ closed.
Choose continuous representatives only through the GNS theorem.  Do not put
\texttt{Continuous} inside the weak-$*$ closed set.

\paragraph{No $\sigma$-compactness shortcut.}
Haar measure on an arbitrary LCA group need not be $\sigma$-finite.  Any use of
$(\Lp1)^*=\Lp\infty$, extraction of almost-everywhere convergent subsequences, or
sequential compactness must be checked for hidden countability assumptions.  The proof
above uses nets, local regularity, and compact testing.  If the available $\Lp1$ duality
theorem is too restrictive, the spectrum proof can use the translation-ratio argument
directly and avoid representing a functional by an $\Lp\infty$ function.

\paragraph{A.e. representatives.}
Convolution of two $\Lp2$ functions and inverse transforms have canonical continuous
representatives.  Give these representatives named definitions and simp lemmas; do not
silently rewrite an $\Lp p$ class as a pointwise function.  Every upgrade from almost
everywhere to everywhere should cite continuity and full support explicitly.

\paragraph{Quotients.}
The compact-lifting lemma used in \cref{thm:dual-quotient} deserves its own general
topological-group declaration: a compact subset of $G/H$ is contained in the image of a
compact subset of $G$.  It avoids any local-section or second-countability assumption.

\paragraph{Noncomputability and normalization.}
Haar measures, Riesz representatives, and the inverse of the evaluation bijection are
noncomputable.  Keep the topological homomorphism \texttt{eval} computable in its formula,
but mark the assembled equivalence and analytic constructions noncomputable.  Store the
dual Haar measure together with its inversion property so that rescaling lemmas can use
uniqueness rather than unfold choices.

\subsection{No-circularity audit}

\begin{longtable}{>{\raggedright\arraybackslash}p{0.30\textwidth}
                  >{\raggedright\arraybackslash}p{0.62\textwidth}}
\toprule
Result & Why it does not presuppose Pontryagin duality \\
\midrule
\endhead
Local compactness of $\dual G$ & Arzel\`a--Ascoli plus the no-small-subgroups property of
$\T$. \\
$\Lp1$ spectrum & Translation in the convolution algebra; no assertion that characters
of $\dual G$ are evaluations. \\
Fourier--Stieltjes uniqueness & Stone--Weierstrass density of transforms of
$\Lp1(G)$; tests only evaluation functions. \\
Bochner & Krein--Milman and the classification of extreme positive coefficients; uses
Fourier--Stieltjes uniqueness only. \\
Dual Haar measure and inversion & Bochner, the multiplier class, Riesz--Markov, and Haar
uniqueness. \\
Injectivity of $\alpha_G$ & The already established inversion formula and a compactly
supported convolution square. \\
Closed image & Topological embedding and the dense-locally-compact-subspace lemma. \\
Surjectivity of $\alpha_G$ & Plancherel supplies a transform supported off the closed
image; asymmetric Fourier--Stieltjes uniqueness forces it to vanish. \\
Character extension & Proved only after biduality, by dualizing the induced quotient map;
it is not an input to the analytic proof. \\
\bottomrule
\end{longtable}

This audit should be mirrored in the Lean import graph.  In particular,
\texttt{Bochner.lean}, \texttt{Inversion.lean}, and \texttt{Plancherel.lean} must not
import \texttt{PontryaginDual/Duality.lean}.

\section{Dependency summary}
\label{sec:dependency-summary}

The logical spine of the proof is:
\begin{enumerate}[label=\arabic*.,leftmargin=2.5em]
\item compact--open characters form an LCA group and joint evaluation is continuous;
\item the $\Lp1$ spectrum is $\dual G$, so Fourier transforms are dense in
$\Czero(\dual G)$ and finite measures are determined by evaluation characters;
\item GNS plus abelian irreducibility identifies extreme normalized positive-type
functions with characters;
\item Banach--Alaoglu, Krein--Milman, and measure uniqueness prove Bochner's theorem;
\item the multiplier construction produces the uniquely normalized Haar measure on
$\dual G$ and Fourier inversion on convolution squares;
\item inversion separates points and recovers the original topology from compact polars;
\item Plancherel produces localized Fourier transforms;
\item a localized transform outside the closed evaluation image contradicts measure
uniqueness, proving surjectivity;
\item biduality then gives annihilator duality, character extension, exactness, and the
categorical anti-equivalence.
\end{enumerate}

Every arrow in this chain points forward.  The proof therefore supplies an
autoformalization order as well as a mathematical proof.

\appendix

\section{Notation-to-Lean dictionary}
\label{app:dictionary}

\begin{longtable}{>{\raggedright\arraybackslash}p{0.23\textwidth}
                  >{\raggedright\arraybackslash}p{0.31\textwidth}
                  >{\raggedright\arraybackslash}p{0.36\textwidth}}
\toprule
Mathematics & Lean target & Comment \\
\midrule
\endhead
$\T$ & \path{Circle} & Subtype of $\C$ of norm one; multiplicative. \\
$\dual G$ & \path{PontryaginDual G} & Definitional equality with
\texttt{G ->\_t* Circle}. \\
$\dual f$ & \path{PontryaginDual.map f} & Precomposition; reverses arrows. \\
$\alpha_G$ & proposed \path{PontryaginDual.eval G} & Continuous homomorphism into the
double dual. \\
$\alpha_G:G\simeq\bidual G$ & proposed \path{PontryaginDual.toDoubleDual G} & A
\path{ContinuousMulEquiv}; name subject to library convention. \\
$W(K,U)$ & compact--open \texttt{MapsTo} set & Use
\path{ContinuousMap.isOpen_setOf_mapsTo}. \\
Haar measure & \path{MeasureTheory.Measure.IsHaarMeasure} / library Haar choice & Keep
the initial measure explicit during Fourier normalization. \\
$f*g$ & group convolution & Be explicit about the selected continuous representative. \\
$f^*$ & involution $x\mapsto\overline{f(x^{-1})}$ & Abelian groups are unimodular. \\
finite complex measure & complex \texttt{VectorMeasure} of finite variation & Positive
special case is an ordinary \texttt{Measure}. \\
$S^\perp$ & proposed \path{PontryaginDual.annihilator S} & Prefer a closed subgroup
when $S$ is a subgroup. \\
$G/H$ & quotient group type & Requires $H$ closed for Hausdorff/LCA instances. \\
\bottomrule
\end{longtable}

\section{Foundational lemmas worth packaging separately}
\label{app:foundation-lemmas}

The proof repeatedly uses the following generic statements.  If any is absent in the
required generality, it should be added outside the Pontryagin namespace.

\begin{enumerate}[label=(A\arabic*)]
\item A dense locally compact subspace of a Hausdorff space is open
(\cref{lem:dense-locally-compact-open}).
\item A locally compact subgroup of a Hausdorff topological group is closed.
This follows by applying (A1) to its closure.
\item Every compact subset of a quotient topological group $G/H$ is contained in the
image of a compact subset of $G$, when $G$ is locally compact.
\item A bounded continuous homomorphism $G\to\C^\times$ has image in $\T$.
\item A positive functional on $\Cc(X)$ is bounded on functions with support in a fixed
compact set.
\item If two continuous functions on a full-support measure space agree almost
everywhere, they agree everywhere.
\item If a continuous density $u$ satisfies $u\mu\geq0$ and $\mu$ has full support, then
$u\geq0$ pointwise.
\item The convolution of two $\Lp2$ functions has a canonical continuous bounded
representative, with the $\Lp2$ norm product bound.
\item Uniform convergence on compact sets of characters, combined with an integrable
tail estimate, gives continuity of the Fourier transform.
\item A positive semidefinite sesquilinear form obeys Cauchy--Schwarz and descends to an
inner product on the quotient by its nullspace.
\end{enumerate}

\section{Bibliographic notes}

The analytic organization follows the classical Haar--Fourier--positive-type route.  The
proof has been rewritten here with conventions and dependency boundaries chosen for
\mathlib.  Useful references for comparison are:

\begin{thebibliography}{9}

\bibitem{HR1}
E. Hewitt and K. A. Ross,
\emph{Abstract Harmonic Analysis, Vol. I},
Springer, second edition, 1979.

\bibitem{HR2}
E. Hewitt and K. A. Ross,
\emph{Abstract Harmonic Analysis, Vol. II},
Springer, 1970.

\bibitem{RV}
D. Ramakrishnan and R. J. Valenza,
\emph{Fourier Analysis on Number Fields},
Graduate Texts in Mathematics 186, Springer, 1999.

\bibitem{Folland}
G. B. Folland,
\emph{A Course in Abstract Harmonic Analysis},
CRC Press, second edition, 2016.

\bibitem{Molteni}
G. Molteni,
\emph{Fourier Analysis in Abelian and Locally Compact Groups},
course notes, Universit\`a degli Studi di Milano.

\bibitem{MathlibPD}
The \mathlib community,
\emph{Mathlib.Topology.Algebra.PontryaginDual},
source and generated documentation for \mathlib~v4.32.0.

\end{thebibliography}

\end{document}
